\documentclass[11pt]{article}
\usepackage[paperwidth=210mm,paperheight=297mm,margin=16mm]{geometry}
\usepackage{fontspec,amsmath,array,multirow,graphicx}
\usepackage{polyglossia}
\setmainlanguage{latin}\setotherlanguages{arabic,greek}
\setmainfont{Linux Libertine O}[Ligatures=TeX]
\newfontfamily\arabicfont{Amiri}[Script=Arabic]
\newfontfamily\greekfont{FreeSerif}[Script=Greek]
\newfontfamily\NallinoSigns{NallinoSigns.otf}[Path=./fonts/]
\newcommand{\AbjadZero}{{\NallinoSigns\symbol{"E000}}}
\newcommand{\AbjadOver}[1]{\leavevmode\vbox{\hrule height .45pt\kern1.2pt\hbox{#1}}}
\newcommand{\Name}[1]{{\addfontfeatures{LetterSpace=12}#1}}
\pagestyle{empty}\setlength{\parindent}{0pt}
\renewcommand{\arraystretch}{1.18}

\usepackage{rotating}
\begin{document}
\clearpage
\noindent\makebox[\textwidth]{\textbf{178}\hfill{\small C. A. NALLINO}\hfill\textbf{}}\par\vspace{-1mm}\noindent\rule{\textwidth}{.4pt}\par\vspace{2mm}
\begin{center}{\small F. 237,v.-239,r.}\end{center}
\begin{center}\footnotesize\begin{tabular}{r|p{30mm}|r@{\hspace{1.5mm}}r|c|r@{\hspace{1.5mm}}r|r@{\hspace{1.5mm}}r|r@{\hspace{1.5mm}}r|r@{\hspace{1.5mm}}r|r@{\hspace{1.5mm}}r||c|}\cline{2-16}
 & \multicolumn{15}{c|}{\rule{0pt}{6mm}\textbf{Status praecipuarum stellarum fixarum anno 1211 a Dhū ’l-qarnayn.}} \\ \cline{2-16}
 & \centering Stellarum nomina. & \multicolumn{2}{c|}{\rotatebox{90}{\parbox{14mm}{\centering\scriptsize Declinatio.}}} & \rotatebox{90}{\parbox{14mm}{\centering\scriptsize Plaga Caeli.}} & \multicolumn{2}{c|}{\rotatebox{90}{\parbox{24mm}{\centering\scriptsize Altitudo meridiana.}}} & \multicolumn{2}{c|}{\rotatebox{90}{\parbox{24mm}{\centering\scriptsize Dimidium commorationis supra Terram.}}} & \multicolumn{2}{c|}{\rotatebox{90}{\parbox{24mm}{\centering\scriptsize Gradus cum quibus culminant.}}} & \multicolumn{2}{c|}{\rotatebox{90}{\parbox{24mm}{\centering\scriptsize Gradus cum quibus oriuntur.}}} & \multicolumn{2}{c|}{\rotatebox{90}{\parbox{24mm}{\centering\scriptsize Gradus cum quibus occidunt.}}} & \rotatebox{90}{\parbox{14mm}{\centering\scriptsize Ordines magnitudinis.}} \\ \hline
1 & \Name{as–Simāk ar–rā-}\newline\hspace*{3mm}\Name{miḥ} (α Bootis). & \raisebox{-1\normalbaselineskip}[0pt][0pt]{25°} & \raisebox{-1\normalbaselineskip}[0pt][0pt]{39′} & \raisebox{-1\normalbaselineskip}[0pt][0pt]{B} & \raisebox{-1\normalbaselineskip}[0pt][0pt]{79°} & \raisebox{-1\normalbaselineskip}[0pt][0pt]{39′} & \raisebox{-1\normalbaselineskip}[0pt][0pt]{110°} & \raisebox{-1\normalbaselineskip}[0pt][0pt]{25′} & \raisebox{-1\normalbaselineskip}[0pt][0pt]{202°} & \raisebox{-1\normalbaselineskip}[0pt][0pt]{48′} & \raisebox{-1\normalbaselineskip}[0pt][0pt]{180°} & \raisebox{-1\normalbaselineskip}[0pt][0pt]{31′} & \raisebox{-1\normalbaselineskip}[0pt][0pt]{239°} & \raisebox{-1\normalbaselineskip}[0pt][0pt]{7′} & \raisebox{-1\normalbaselineskip}[0pt][0pt]{1} \\
2 & \Name{an–Nasr al–wāqiʿ}\newline\hspace*{3mm}(α Lyrae). & \raisebox{-1\normalbaselineskip}[0pt][0pt]{38} & \raisebox{-1\normalbaselineskip}[0pt][0pt]{25} & \raisebox{-1\normalbaselineskip}[0pt][0pt]{B} & \raisebox{-1\normalbaselineskip}[0pt][0pt]{92} & \raisebox{-1\normalbaselineskip}[0pt][0pt]{25} & \raisebox{-1\normalbaselineskip}[0pt][0pt]{125} & \raisebox{-1\normalbaselineskip}[0pt][0pt]{8} & \raisebox{-1\normalbaselineskip}[0pt][0pt]{269} & \raisebox{-1\normalbaselineskip}[0pt][0pt]{20} & \raisebox{-1\normalbaselineskip}[0pt][0pt]{224} & \raisebox{-1\normalbaselineskip}[0pt][0pt]{25} & \raisebox{-1\normalbaselineskip}[0pt][0pt]{314} & \raisebox{-1\normalbaselineskip}[0pt][0pt]{25} & \raisebox{-1\normalbaselineskip}[0pt][0pt]{1} \\
3 & \Name{al–ʿAnz} sive \Name{al–ʿAy-}\newline\hspace*{3mm}\Name{yūq} (α Aurigae). & \raisebox{-1\normalbaselineskip}[0pt][0pt]{43} & \raisebox{-1\normalbaselineskip}[0pt][0pt]{10} & \raisebox{-1\normalbaselineskip}[0pt][0pt]{B} & \raisebox{-1\normalbaselineskip}[0pt][0pt]{97} & \raisebox{-1\normalbaselineskip}[0pt][0pt]{10} & \raisebox{-1\normalbaselineskip}[0pt][0pt]{132} & \raisebox{-1\normalbaselineskip}[0pt][0pt]{58} & \raisebox{-1\normalbaselineskip}[0pt][0pt]{61} & \raisebox{-1\normalbaselineskip}[0pt][0pt]{34} & \raisebox{-1\normalbaselineskip}[0pt][0pt]{25} & \raisebox{-1\normalbaselineskip}[0pt][0pt]{45} & \raisebox{-1\normalbaselineskip}[0pt][0pt]{84} & \raisebox{-1\normalbaselineskip}[0pt][0pt]{28} & \raisebox{-1\normalbaselineskip}[0pt][0pt]{1} \\
4 & \Name{ad–Dabarān} (α Tau-\newline\hspace*{3mm}ri). & \raisebox{-1\normalbaselineskip}[0pt][0pt]{13} & \raisebox{-1\normalbaselineskip}[0pt][0pt]{19} & \raisebox{-1\normalbaselineskip}[0pt][0pt]{B} & \raisebox{-1\normalbaselineskip}[0pt][0pt]{67} & \raisebox{-1\normalbaselineskip}[0pt][0pt]{19} & \raisebox{-1\normalbaselineskip}[0pt][0pt]{99} & \raisebox{-1\normalbaselineskip}[0pt][0pt]{54} & \raisebox{-1\normalbaselineskip}[0pt][0pt]{55} & \raisebox{-1\normalbaselineskip}[0pt][0pt]{25} & \raisebox{-1\normalbaselineskip}[0pt][0pt]{61} & \raisebox{-1\normalbaselineskip}[0pt][0pt]{4} & \raisebox{-1\normalbaselineskip}[0pt][0pt]{51} & \raisebox{-1\normalbaselineskip}[0pt][0pt]{32} & \raisebox{-1\normalbaselineskip}[0pt][0pt]{1} \\
5 & \Name{Qalb al–asad} (α\newline\hspace*{3mm}Leonis). & \raisebox{-1\normalbaselineskip}[0pt][0pt]{16} & \raisebox{-1\normalbaselineskip}[0pt][0pt]{47} & \raisebox{-1\normalbaselineskip}[0pt][0pt]{B} & \raisebox{-1\normalbaselineskip}[0pt][0pt]{70} & \raisebox{-1\normalbaselineskip}[0pt][0pt]{47} & \raisebox{-1\normalbaselineskip}[0pt][0pt]{102} & \raisebox{-1\normalbaselineskip}[0pt][0pt]{50} & \raisebox{-1\normalbaselineskip}[0pt][0pt]{134} & \raisebox{-1\normalbaselineskip}[0pt][0pt]{20} & \raisebox{-1\normalbaselineskip}[0pt][0pt]{134} & \raisebox{-1\normalbaselineskip}[0pt][0pt]{5} & \raisebox{-1\normalbaselineskip}[0pt][0pt]{134} & \raisebox{-1\normalbaselineskip}[0pt][0pt]{43} & \raisebox{-1\normalbaselineskip}[0pt][0pt]{1} \\
6 & \Name{Dhanab al–asad}\newline\hspace*{3mm}(β Leonis). & \raisebox{-1\normalbaselineskip}[0pt][0pt]{19} & \raisebox{-1\normalbaselineskip}[0pt][0pt]{36} & \raisebox{-1\normalbaselineskip}[0pt][0pt]{B} & \raisebox{-1\normalbaselineskip}[0pt][0pt]{73} & \raisebox{-1\normalbaselineskip}[0pt][0pt]{36} & \raisebox{-1\normalbaselineskip}[0pt][0pt]{105} & \raisebox{-1\normalbaselineskip}[0pt][0pt]{1} & \raisebox{-1\normalbaselineskip}[0pt][0pt]{160} & \raisebox{-1\normalbaselineskip}[0pt][0pt]{52} & \raisebox{-1\normalbaselineskip}[0pt][0pt]{153} & \raisebox{-1\normalbaselineskip}[0pt][0pt]{5} & \raisebox{-1\normalbaselineskip}[0pt][0pt]{175} & \raisebox{-1\normalbaselineskip}[0pt][0pt]{48} & \raisebox{-1\normalbaselineskip}[0pt][0pt]{1} \\
7 & Postrema cursus a-\newline\hspace*{3mm}quae (α Piscis au-\newline\hspace*{3mm}stralis). & \raisebox{-2\normalbaselineskip}[0pt][0pt]{\textit{33}} & \raisebox{-2\normalbaselineskip}[0pt][0pt]{\textit{45}} & \raisebox{-2\normalbaselineskip}[0pt][0pt]{A} & \raisebox{-2\normalbaselineskip}[0pt][0pt]{\textit{20}} & \raisebox{-2\normalbaselineskip}[0pt][0pt]{\textit{15}} & \raisebox{-2\normalbaselineskip}[0pt][0pt]{63} & \raisebox{-2\normalbaselineskip}[0pt][0pt]{50} & \raisebox{-2\normalbaselineskip}[0pt][0pt]{325} & \raisebox{-2\normalbaselineskip}[0pt][0pt]{25} & \raisebox{-2\normalbaselineskip}[0pt][0pt]{350} & \raisebox{-2\normalbaselineskip}[0pt][0pt]{15} & \raisebox{-2\normalbaselineskip}[0pt][0pt]{312} & \raisebox{-2\normalbaselineskip}[0pt][0pt]{6} & \raisebox{-2\normalbaselineskip}[0pt][0pt]{1} \\
\cline{2-16}\end{tabular}\end{center}
{\small
\par\hangindent=12mm\hangafter=1 1. Cum \textit{P} (1) consentiunt δ et \textit{h}. — \textit{A} potius pro anno 1216 valeret; tamen pro anno 1211 error nonnisi + 4′ est. Datis \textit{P} et \textit{A} bene dati \textit{B} et \textit{C} deducuntur.
\par\hangindent=12mm\hangafter=1 2. In \textit{h} cod. \textarabic{ضو} 96° pro \textarabic{ضٮ} = \textarabic{ضب} 92°; in \textit{B} cod. \textarabic{ٯكد} = \textarabic{قكد} 124° pro \textarabic{ركد} 224°. — \textit{P} differt − 3′ ab eo qui cum δ et \textit{h} conveniret. Optime \textit{A}, cum qua et cum \textit{P} conveniunt \textit{B} et \textit{C}.
\par\hangindent=12mm\hangafter=1 3. In \textit{h} cod. \textarabic{صز} 67° pro \textarabic{ضز} 97°; in \textit{C} cod. \textarabic{قز} 107° pro \textarabic{فد} 84°. — \textit{P}, δ, \textit{h} inter se convenientes. \textit{A} differt + 2′ vel + 3′ a vera quantitate pro 1211.
\par\hangindent=12mm\hangafter=1 4. In \textit{A} cod. \textarabic{له} 35° pro \textarabic{نه} 55°; in \textit{B} \textarabic{ى} 10′ pro \textarabic{د} 4′; in \textit{C} \textarabic{رنا} 251° pro \textarabic{نا} 51°.
\par\hangindent=12mm\hangafter=1 5. Ex δ et \textit{h} proveniret \textit{P} 102° 40′ loco codicis 102° 50′; sed hoc numero auctor usus est ad \textit{B} et \textit{C} supputandos. Tam δ quam \textit{A} concordant cum catalogo ad annum 1211 reducto, quod sane perraro in his tabulis evenit.
\par\hangindent=12mm\hangafter=1 6. \textit{P}, δ, \textit{h} satis inter se convenientes; \textit{A} quantitate 3′ nimis parva.
\par\hangindent=12mm\hangafter=1 7. Plaga cod. borealis. — Magnus error in δ et \textit{h}. In \textit{P} cod. \textarabic{سج} 303° pro \textarabic{صج} 63°; in \textit{B} cod. \textarabic{سز م} 307° 40′ pro \textarabic{سن يه} 350° 15′. Quibus correctis, data \textit{A} (quae − 2° differt ab anno 1211) bene \textit{B} et \textit{C} deducuntur.
\par}
\par\vspace{2mm}\begin{center}\rule{40mm}{.4pt}\end{center}{\small (1) Adnotationes omnes ad hanc tabulam nituntur longissimis taediosissimisque supputationibus quibus \Name{Schiaparelli} multos codicis errores, quantum fieri poterat, emendare conatus est. Multa incerta semper esse patet quoque ex animadversionibus generalibus sub fine huius tomi impressis. — Littera δ declinationem, \textit{h} altitudinem, \textit{P} arcum semidiurnum (dimidium commorationis supra Terram), \textit{A} culminationem, \textit{B} ortum, \textit{C} occasum indico. — De magnitudinibus, praeter animadversiones generales, cfr. p. 183, adn. 1.}
\clearpage
\noindent\makebox[\textwidth]{\textbf{}\hfill{\small AL-BATTANI OPUS ASTRONOMICUM}\hfill\textbf{179}}\par\vspace{-1mm}\noindent\rule{\textwidth}{.4pt}\par\vspace{2mm}
\begin{center}\footnotesize\begin{tabular}{r|p{30mm}|r@{\hspace{1.5mm}}r|c|r@{\hspace{1.5mm}}r|r@{\hspace{1.5mm}}r|r@{\hspace{1.5mm}}r|r@{\hspace{1.5mm}}r|r@{\hspace{1.5mm}}r||c|}\cline{2-16}
 & \centering Stellarum nomina. & \multicolumn{2}{c|}{\rotatebox{90}{\parbox{14mm}{\centering\scriptsize Declinatio.}}} & \rotatebox{90}{\parbox{14mm}{\centering\scriptsize Plaga Caeli.}} & \multicolumn{2}{c|}{\rotatebox{90}{\parbox{24mm}{\centering\scriptsize Altitudo meridiana.}}} & \multicolumn{2}{c|}{\rotatebox{90}{\parbox{24mm}{\centering\scriptsize Dimidium commorationis supra Terram.}}} & \multicolumn{2}{c|}{\rotatebox{90}{\parbox{24mm}{\centering\scriptsize Gradus cum quibus culminant.}}} & \multicolumn{2}{c|}{\rotatebox{90}{\parbox{24mm}{\centering\scriptsize Gradus cum quibus oriuntur.}}} & \multicolumn{2}{c|}{\rotatebox{90}{\parbox{24mm}{\centering\scriptsize Gradus cum quibus occidunt.}}} & \rotatebox{90}{\parbox{14mm}{\centering\scriptsize Ordines magnitudinis.}} \\ \hline
8 & \Name{as–Simāk al–aʿzal}\newline\hspace*{3mm}(α Virginis). & \raisebox{-1\normalbaselineskip}[0pt][0pt]{4°} & \raisebox{-1\normalbaselineskip}[0pt][0pt]{49′} & \raisebox{-1\normalbaselineskip}[0pt][0pt]{A} & \raisebox{-1\normalbaselineskip}[0pt][0pt]{49°} & \raisebox{-1\normalbaselineskip}[0pt][0pt]{11′} & \raisebox{-1\normalbaselineskip}[0pt][0pt]{86°} & \raisebox{-1\normalbaselineskip}[0pt][0pt]{31′} & \raisebox{-1\normalbaselineskip}[0pt][0pt]{186°} & \raisebox{-1\normalbaselineskip}[0pt][0pt]{43′} & \raisebox{-1\normalbaselineskip}[0pt][0pt]{188°} & \raisebox{-1\normalbaselineskip}[0pt][0pt]{0′} & \raisebox{-1\normalbaselineskip}[0pt][0pt]{184°} & \raisebox{-1\normalbaselineskip}[0pt][0pt]{20′} & \raisebox{-1\normalbaselineskip}[0pt][0pt]{1} \\
9 & Initium caudae Sa-\newline\hspace*{3mm}gittarii (ω Sagitta-\newline\hspace*{3mm}rii). & \raisebox{-2\normalbaselineskip}[0pt][0pt]{30} & \raisebox{-2\normalbaselineskip}[0pt][0pt]{0} & \raisebox{-2\normalbaselineskip}[0pt][0pt]{A} & \raisebox{-2\normalbaselineskip}[0pt][0pt]{24} & \raisebox{-2\normalbaselineskip}[0pt][0pt]{0} & \raisebox{-2\normalbaselineskip}[0pt][0pt]{65} & \raisebox{-2\normalbaselineskip}[0pt][0pt]{12} & \raisebox{-2\normalbaselineskip}[0pt][0pt]{274} & \raisebox{-2\normalbaselineskip}[0pt][0pt]{30} & \raisebox{-2\normalbaselineskip}[0pt][0pt]{280} & \raisebox{-2\normalbaselineskip}[0pt][0pt]{40} & \raisebox{-2\normalbaselineskip}[0pt][0pt]{268} & \raisebox{-2\normalbaselineskip}[0pt][0pt]{43} & \raisebox{-2\normalbaselineskip}[0pt][0pt]{1} \\
10 & \Name{Mankib al–ǵawzā’}\newline\hspace*{3mm}(α Orionis). & \raisebox{-1\normalbaselineskip}[0pt][0pt]{5} & \raisebox{-1\normalbaselineskip}[0pt][0pt]{30} & \raisebox{-1\normalbaselineskip}[0pt][0pt]{B} & \raisebox{-1\normalbaselineskip}[0pt][0pt]{59} & \raisebox{-1\normalbaselineskip}[0pt][0pt]{30} & \raisebox{-1\normalbaselineskip}[0pt][0pt]{94} & \raisebox{-1\normalbaselineskip}[0pt][0pt]{0} & \raisebox{-1\normalbaselineskip}[0pt][0pt]{76} & \raisebox{-1\normalbaselineskip}[0pt][0pt]{25} & \raisebox{-1\normalbaselineskip}[0pt][0pt]{89} & \raisebox{-1\normalbaselineskip}[0pt][0pt]{45} & \raisebox{-1\normalbaselineskip}[0pt][0pt]{64} & \raisebox{-1\normalbaselineskip}[0pt][0pt]{50} & \raisebox{-1\normalbaselineskip}[0pt][0pt]{1} \\
11 & \Name{Riǵl al-ǵawzā’} (β\newline\hspace*{3mm}Orionis). & \raisebox{-1\normalbaselineskip}[0pt][0pt]{10} & \raisebox{-1\normalbaselineskip}[0pt][0pt]{20} & \raisebox{-1\normalbaselineskip}[0pt][0pt]{A} & \raisebox{-1\normalbaselineskip}[0pt][0pt]{43} & \raisebox{-1\normalbaselineskip}[0pt][0pt]{40} & \raisebox{-1\normalbaselineskip}[0pt][0pt]{82} & \raisebox{-1\normalbaselineskip}[0pt][0pt]{21} & \raisebox{-1\normalbaselineskip}[0pt][0pt]{69} & \raisebox{-1\normalbaselineskip}[0pt][0pt]{56} & \raisebox{-1\normalbaselineskip}[0pt][0pt]{94} & \raisebox{-1\normalbaselineskip}[0pt][0pt]{0} & \raisebox{-1\normalbaselineskip}[0pt][0pt]{49} & \raisebox{-1\normalbaselineskip}[0pt][0pt]{40} & \raisebox{-1\normalbaselineskip}[0pt][0pt]{1} \\
12 & Postrema stellarum\newline\hspace*{3mm}Fluvii (θ Eridani). & \raisebox{-1\normalbaselineskip}[0pt][0pt]{\textit{43}} & \raisebox{-1\normalbaselineskip}[0pt][0pt]{\textit{45}} & \raisebox{-1\normalbaselineskip}[0pt][0pt]{A} & \raisebox{-1\normalbaselineskip}[0pt][0pt]{\textit{10}} & \raisebox{-1\normalbaselineskip}[0pt][0pt]{\textit{15}} & \raisebox{-1\normalbaselineskip}[0pt][0pt]{45} & \raisebox{-1\normalbaselineskip}[0pt][0pt]{48} & \raisebox{-1\normalbaselineskip}[0pt][0pt]{39} & \raisebox{-1\normalbaselineskip}[0pt][0pt]{3} & \raisebox{-1\normalbaselineskip}[0pt][0pt]{98} & \raisebox{-1\normalbaselineskip}[0pt][0pt]{20} & \raisebox{-1\normalbaselineskip}[0pt][0pt]{353} & \raisebox{-1\normalbaselineskip}[0pt][0pt]{43} & \raisebox{-1\normalbaselineskip}[0pt][0pt]{1} \\
13 & \Name{ash–Shiʿrà ’l–ya-}\newline\hspace*{3mm}\Name{māniyyah} (α Ca-\newline\hspace*{3mm}nis maioris). & \raisebox{-2\normalbaselineskip}[0pt][0pt]{15} & \raisebox{-2\normalbaselineskip}[0pt][0pt]{51} & \raisebox{-2\normalbaselineskip}[0pt][0pt]{A} & \raisebox{-2\normalbaselineskip}[0pt][0pt]{38} & \raisebox{-2\normalbaselineskip}[0pt][0pt]{9} & \raisebox{-2\normalbaselineskip}[0pt][0pt]{78} & \raisebox{-2\normalbaselineskip}[0pt][0pt]{5} & \raisebox{-2\normalbaselineskip}[0pt][0pt]{89} & \raisebox{-2\normalbaselineskip}[0pt][0pt]{22} & \raisebox{-2\normalbaselineskip}[0pt][0pt]{115} & \raisebox{-2\normalbaselineskip}[0pt][0pt]{33} & \raisebox{-2\normalbaselineskip}[0pt][0pt]{63} & \raisebox{-2\normalbaselineskip}[0pt][0pt]{18} & \raisebox{-2\normalbaselineskip}[0pt][0pt]{1} \\
14 & \Name{ash–Shiʿrà ’sh-}\newline\hspace*{3mm}\Name{–sha’miyyah} (α\newline\hspace*{3mm}Canis minoris). & \raisebox{-2\normalbaselineskip}[0pt][0pt]{6} & \raisebox{-2\normalbaselineskip}[0pt][0pt]{25} & \raisebox{-2\normalbaselineskip}[0pt][0pt]{B} & \raisebox{-2\normalbaselineskip}[0pt][0pt]{60} & \raisebox{-2\normalbaselineskip}[0pt][0pt]{25} & \raisebox{-2\normalbaselineskip}[0pt][0pt]{94} & \raisebox{-2\normalbaselineskip}[0pt][0pt]{41} & \raisebox{-2\normalbaselineskip}[0pt][0pt]{99} & \raisebox{-2\normalbaselineskip}[0pt][0pt]{25} & \raisebox{-2\normalbaselineskip}[0pt][0pt]{110} & \raisebox{-2\normalbaselineskip}[0pt][0pt]{53} & \raisebox{-2\normalbaselineskip}[0pt][0pt]{86} & \raisebox{-2\normalbaselineskip}[0pt][0pt]{46} & \raisebox{-2\normalbaselineskip}[0pt][0pt]{1} \\
15 & \Name{Suhayl al–yamānī}\newline\hspace*{3mm}(Canopus). & \raisebox{-1\normalbaselineskip}[0pt][0pt]{51} & \raisebox{-1\normalbaselineskip}[0pt][0pt]{25} & \raisebox{-1\normalbaselineskip}[0pt][0pt]{A} & \raisebox{-1\normalbaselineskip}[0pt][0pt]{2} & \raisebox{-1\normalbaselineskip}[0pt][0pt]{35} & \raisebox{-1\normalbaselineskip}[0pt][0pt]{24} & \raisebox{-1\normalbaselineskip}[0pt][0pt]{26} & \raisebox{-1\normalbaselineskip}[0pt][0pt]{89} & \raisebox{-1\normalbaselineskip}[0pt][0pt]{30} & \raisebox{-1\normalbaselineskip}[0pt][0pt]{159} & \raisebox{-1\normalbaselineskip}[0pt][0pt]{20} & \raisebox{-1\normalbaselineskip}[0pt][0pt]{19} & \raisebox{-1\normalbaselineskip}[0pt][0pt]{44} & \raisebox{-1\normalbaselineskip}[0pt][0pt]{1} \\
16 & \textit{Riǵl Qinṭāwuros}\newline\hspace*{3mm}« \textit{Pes Centauri} »\newline\hspace*{3mm}(β Columbae). & \raisebox{-2\normalbaselineskip}[0pt][0pt]{37} & \raisebox{-2\normalbaselineskip}[0pt][0pt]{3} & \raisebox{-2\normalbaselineskip}[0pt][0pt]{A} & \raisebox{-2\normalbaselineskip}[0pt][0pt]{16} & \raisebox{-2\normalbaselineskip}[0pt][0pt]{57} & \raisebox{-2\normalbaselineskip}[0pt][0pt]{56} & \raisebox{-2\normalbaselineskip}[0pt][0pt]{44} & \raisebox{-2\normalbaselineskip}[0pt][0pt]{78} & \raisebox{-2\normalbaselineskip}[0pt][0pt]{50} & \raisebox{-2\normalbaselineskip}[0pt][0pt]{123} & \raisebox{-2\normalbaselineskip}[0pt][0pt]{36} & \raisebox{-2\normalbaselineskip}[0pt][0pt]{36} & \raisebox{-2\normalbaselineskip}[0pt][0pt]{39} & \raisebox{-2\normalbaselineskip}[0pt][0pt]{2} \\
\cline{2-16}\end{tabular}\end{center}
{\small
\par\hangindent=12mm\hangafter=1 8. Plaga cod. borealis. In \textit{h} cod. \textarabic{ما} 41′ pro \textarabic{ٮا} = \textarabic{يا} 11′; in \textit{P} cod. \textarabic{قو} 106° pro \textarabic{فو} 86°; in \textit{B} cod. \textarabic{قعج} 173° pro \textarabic{قڡح} = \textarabic{قفح} 188°; in \textit{C} cod. \textarabic{قعد} 174° pro \textarabic{قڡد} = \textarabic{قفد} 184°. — \textit{P} et δ bene inter se convenientes; a \textit{P} et \textit{A} optime \textit{B} et \textit{C} proveniunt.
\par\hangindent=12mm\hangafter=1 9. Plaga cod. borealis. E catalogo suo fixarum al–Battānī long. 274° 0′, lat. 6° 30′ sumpsit (cfr. p. 163, nr. 16), cum quibus omnia optime congruunt.
\par\hangindent=12mm\hangafter=1 10. In \textit{B} cod. \textarabic{قط} 109° pro \textarabic{فط} 89°. \textit{P}, \textit{h}, δ optime inter se conveniunt; item reliqua.
\par\hangindent=12mm\hangafter=1 11. Plaga cod. borealis. In \textit{P} cod. \textarabic{قب} 102° pro \textarabic{فب} 82°; in \textit{C} cod. \textarabic{مد \AbjadZero{}} 44° 0′ pro \textarabic{مط م} 49° 40′. — \textit{P} satis convenit cum δ et \textit{h}, sed hae et \textit{A} discrepant a catalogo ad annum 1211 reducto. Tamen numeris datis exacte \textit{B} et \textit{C} deducuntur.
\par\hangindent=12mm\hangafter=1 12. De hac stella cfr. p. 170, nr. 2. — Plaga cod. borealis. A catalogo fixarum deducerentur δ 43° 32′, \textit{P} 46° 21′, \textit{A} 38° 57′. Sed a codicis \textit{P} et \textit{A} exacte deducuntur \textit{B} et \textit{C}. Cum \textit{P} cod. conveniret δ 43° 49′ pro 43° 45′; error in scriptura maghrebina facillimus.
\par\hangindent=12mm\hangafter=1 13. Plaga cod. borealis. In \textit{P} cod. \textarabic{ن} 50′ pro \textarabic{ه} 5′; in \textit{A} \textarabic{قط} 109° pro \textarabic{فط} 89°; in \textit{C} \textarabic{سنج} 353° pro \textarabic{صج} 63° (in archetypo, iuxta scripturam orientalem, \textarabic{سج} 63°). Reliqua optime.
\par\hangindent=12mm\hangafter=1 14. In \textit{C} cod. \textarabic{قو} 106° pro \textarabic{فو} 86°.
\par\hangindent=12mm\hangafter=1 15. In \textit{A} cod. \textarabic{قط} 109° pro \textarabic{فط} 89°; in \textit{B} cod. \textarabic{قيط} 119° pro \textarabic{قنط} 159°; in \textit{C} \textarabic{نط} 59° pro \textarabic{يط} 19°. Quibus correctis, \textit{A}, \textit{B}, \textit{C} inter se conveniunt. Cum \textit{P} et cum anno 1211 satis bene consentiunt δ et \textit{h}.
\par\hangindent=12mm\hangafter=1 16. Stellae nomen falsum; non enim est α Centauri, sed β Columbae (nr. 9 constellationis Canis maioris in catalogo Albateniano, p. 172). Vero suo nomine rursus indicatur infra nr. 60; elementa autem ad
\par}
\clearpage
\noindent\makebox[\textwidth]{\textbf{180}\hfill{\small C. A. NALLINO}\hfill\textbf{}}\par\vspace{-1mm}\noindent\rule{\textwidth}{.4pt}\par\vspace{2mm}
\begin{center}\footnotesize\begin{tabular}{r|p{30mm}|r@{\hspace{1.5mm}}r|c|r@{\hspace{1.5mm}}r|r@{\hspace{1.5mm}}r|r@{\hspace{1.5mm}}r|r@{\hspace{1.5mm}}r|r@{\hspace{1.5mm}}r||c|}\cline{2-16}
 & \centering Stellarum nomina. & \multicolumn{2}{c|}{\rotatebox{90}{\parbox{14mm}{\centering\scriptsize Declinatio.}}} & \rotatebox{90}{\parbox{14mm}{\centering\scriptsize Plaga Caeli.}} & \multicolumn{2}{c|}{\rotatebox{90}{\parbox{24mm}{\centering\scriptsize Altitudo meridiana.}}} & \multicolumn{2}{c|}{\rotatebox{90}{\parbox{24mm}{\centering\scriptsize Dimidium commorationis supra Terram.}}} & \multicolumn{2}{c|}{\rotatebox{90}{\parbox{24mm}{\centering\scriptsize Gradus cum quibus culminant.}}} & \multicolumn{2}{c|}{\rotatebox{90}{\parbox{24mm}{\centering\scriptsize Gradus cum quibus oriuntur.}}} & \multicolumn{2}{c|}{\rotatebox{90}{\parbox{24mm}{\centering\scriptsize Gradus cum quibus occidunt.}}} & \rotatebox{90}{\parbox{14mm}{\centering\scriptsize Ordines magnitudinis.}} \\ \hline
17 & Latus Ursae minoris\newline\hspace*{3mm}(β Ursae minoris). & \raisebox{-1\normalbaselineskip}[0pt][0pt]{\textit{91°}} & \raisebox{-1\normalbaselineskip}[0pt][0pt]{\textit{17′}} & \raisebox{-1\normalbaselineskip}[0pt][0pt]{B} & \raisebox{-1\normalbaselineskip}[0pt][0pt]{\textit{145°}} & \raisebox{-1\normalbaselineskip}[0pt][0pt]{\textit{17′}} & \multicolumn{2}{c|}{\multirow{8}{*}{\rotatebox{90}{\parbox{48mm}{\centering\scriptsize Hae octo stellae supra Terram semper apparent;\\nec ortum nec occasum habent, et in suis revolutio-}}}} & \raisebox{-1\normalbaselineskip}[0pt][0pt]{\textit{130°}} & \raisebox{-1\normalbaselineskip}[0pt][0pt]{\textit{0′}} & \multicolumn{4}{c|}{\multirow{8}{*}{\rotatebox{90}{\parbox{48mm}{\centering\scriptsize nibus lineam medii Caeli bis attingunt, altera vice su-\\pra polum cum gradibus in hac tabula descriptis, al-\\tera vice sub polo.}}}} & \raisebox{-1\normalbaselineskip}[0pt][0pt]{2} \\
18 & Eius latus etiam (γ\newline\hspace*{3mm}Ursae minoris). & \raisebox{-1\normalbaselineskip}[0pt][0pt]{\textit{91}} & \raisebox{-1\normalbaselineskip}[0pt][0pt]{\textit{15}} & \raisebox{-1\normalbaselineskip}[0pt][0pt]{B} & \raisebox{-1\normalbaselineskip}[0pt][0pt]{\textit{145}} & \raisebox{-1\normalbaselineskip}[0pt][0pt]{\textit{15}} & \multicolumn{2}{c|}{} & \raisebox{-1\normalbaselineskip}[0pt][0pt]{\textit{142}} & \raisebox{-1\normalbaselineskip}[0pt][0pt]{\textit{24}} & \multicolumn{4}{c|}{} & \raisebox{-1\normalbaselineskip}[0pt][0pt]{2} \\
19 & Dorsus Ursae maio-\newline\hspace*{3mm}ris (α Ursae maio-\newline\hspace*{3mm}ris). & \raisebox{-2\normalbaselineskip}[0pt][0pt]{66} & \raisebox{-2\normalbaselineskip}[0pt][0pt]{20} & \raisebox{-2\normalbaselineskip}[0pt][0pt]{B} & \raisebox{-2\normalbaselineskip}[0pt][0pt]{120} & \raisebox{-2\normalbaselineskip}[0pt][0pt]{20} & \multicolumn{2}{c|}{} & \raisebox{-2\normalbaselineskip}[0pt][0pt]{142} & \raisebox{-2\normalbaselineskip}[0pt][0pt]{0} & \multicolumn{4}{c|}{} & \raisebox{-2\normalbaselineskip}[0pt][0pt]{2} \\
20 & Tenuior pars ventris\newline\hspace*{3mm}eius (β Ursae ma-\newline\hspace*{3mm}ioris). & \raisebox{-2\normalbaselineskip}[0pt][0pt]{60} & \raisebox{-2\normalbaselineskip}[0pt][0pt]{30} & \raisebox{-2\normalbaselineskip}[0pt][0pt]{B} & \raisebox{-2\normalbaselineskip}[0pt][0pt]{114} & \raisebox{-2\normalbaselineskip}[0pt][0pt]{30} & \multicolumn{2}{c|}{} & \raisebox{-2\normalbaselineskip}[0pt][0pt]{143} & \raisebox{-2\normalbaselineskip}[0pt][0pt]{10} & \multicolumn{4}{c|}{} & \raisebox{-2\normalbaselineskip}[0pt][0pt]{2} \\
21 & Eius femur (γ Ur-\newline\hspace*{3mm}sae maioris). & \raisebox{-1\normalbaselineskip}[0pt][0pt]{57} & \raisebox{-1\normalbaselineskip}[0pt][0pt]{45} & \raisebox{-1\normalbaselineskip}[0pt][0pt]{B} & \raisebox{-1\normalbaselineskip}[0pt][0pt]{111} & \raisebox{-1\normalbaselineskip}[0pt][0pt]{45} & \multicolumn{2}{c|}{} & \raisebox{-1\normalbaselineskip}[0pt][0pt]{\textit{178}} & \raisebox{-1\normalbaselineskip}[0pt][0pt]{\textit{35}} & \multicolumn{4}{c|}{} & \raisebox{-1\normalbaselineskip}[0pt][0pt]{2} \\
22 & Eius cauda (ε Ursae\newline\hspace*{3mm}maioris). & \raisebox{-1\normalbaselineskip}[0pt][0pt]{60} & \raisebox{-1\normalbaselineskip}[0pt][0pt]{27} & \raisebox{-1\normalbaselineskip}[0pt][0pt]{B} & \raisebox{-1\normalbaselineskip}[0pt][0pt]{114} & \raisebox{-1\normalbaselineskip}[0pt][0pt]{26} & \multicolumn{2}{c|}{} & \raisebox{-1\normalbaselineskip}[0pt][0pt]{\textit{168}} & \raisebox{-1\normalbaselineskip}[0pt][0pt]{\textit{10}} & \multicolumn{4}{c|}{} & \raisebox{-1\normalbaselineskip}[0pt][0pt]{2} \\
23 & Eius cauda etiam (ζ\newline\hspace*{3mm}Ursae maioris). & \raisebox{-1\normalbaselineskip}[0pt][0pt]{59} & \raisebox{-1\normalbaselineskip}[0pt][0pt]{40} & \raisebox{-1\normalbaselineskip}[0pt][0pt]{B} & \raisebox{-1\normalbaselineskip}[0pt][0pt]{113} & \raisebox{-1\normalbaselineskip}[0pt][0pt]{40} & \multicolumn{2}{c|}{} & \raisebox{-1\normalbaselineskip}[0pt][0pt]{187} & \raisebox{-1\normalbaselineskip}[0pt][0pt]{20} & \multicolumn{4}{c|}{} & \raisebox{-1\normalbaselineskip}[0pt][0pt]{2} \\
24 & Eius cauda etiam (η\newline\hspace*{3mm}Ursae maioris). & \raisebox{-1\normalbaselineskip}[0pt][0pt]{54} & \raisebox{-1\normalbaselineskip}[0pt][0pt]{15} & \raisebox{-1\normalbaselineskip}[0pt][0pt]{B} & \raisebox{-1\normalbaselineskip}[0pt][0pt]{108} & \raisebox{-1\normalbaselineskip}[0pt][0pt]{15} & \multicolumn{2}{c|}{} & \raisebox{-1\normalbaselineskip}[0pt][0pt]{196} & \raisebox{-1\normalbaselineskip}[0pt][0pt]{0} & \multicolumn{4}{c|}{} & \raisebox{-1\normalbaselineskip}[0pt][0pt]{2} \\
25 & Lucida Coronae bo-\newline\hspace*{3mm}realis (α Coronae\newline\hspace*{3mm}bor.). & \raisebox{-2\normalbaselineskip}[0pt][0pt]{31} & \raisebox{-2\normalbaselineskip}[0pt][0pt]{46} & \raisebox{-2\normalbaselineskip}[0pt][0pt]{B} & \raisebox{-2\normalbaselineskip}[0pt][0pt]{85} & \raisebox{-2\normalbaselineskip}[0pt][0pt]{46} & \raisebox{-2\normalbaselineskip}[0pt][0pt]{\textit{134°}} & \raisebox{-2\normalbaselineskip}[0pt][0pt]{\textit{18′}} & \raisebox{-2\normalbaselineskip}[0pt][0pt]{224} & \raisebox{-2\normalbaselineskip}[0pt][0pt]{11} & \raisebox{-2\normalbaselineskip}[0pt][0pt]{192°} & \raisebox{-2\normalbaselineskip}[0pt][0pt]{32′} & \raisebox{-2\normalbaselineskip}[0pt][0pt]{267°} & \raisebox{-2\normalbaselineskip}[0pt][0pt]{6′} & \raisebox{-2\normalbaselineskip}[0pt][0pt]{2} \\
26 & Cauda Cygni (α Cy-\newline\hspace*{3mm}gni). & \raisebox{-1\normalbaselineskip}[0pt][0pt]{42} & \raisebox{-1\normalbaselineskip}[0pt][0pt]{21} & \raisebox{-1\normalbaselineskip}[0pt][0pt]{B} & \raisebox{-1\normalbaselineskip}[0pt][0pt]{96} & \raisebox{-1\normalbaselineskip}[0pt][0pt]{21} & \raisebox{-1\normalbaselineskip}[0pt][0pt]{131} & \raisebox{-1\normalbaselineskip}[0pt][0pt]{27} & \raisebox{-1\normalbaselineskip}[0pt][0pt]{298} & \raisebox{-1\normalbaselineskip}[0pt][0pt]{54} & \raisebox{-1\normalbaselineskip}[0pt][0pt]{245} & \raisebox{-1\normalbaselineskip}[0pt][0pt]{7} & \raisebox{-1\normalbaselineskip}[0pt][0pt]{345} & \raisebox{-1\normalbaselineskip}[0pt][0pt]{31} & \raisebox{-1\normalbaselineskip}[0pt][0pt]{2} \\
27 & Rostrum Cygni (β Cy-\newline\hspace*{3mm}gni). & \raisebox{-1\normalbaselineskip}[0pt][0pt]{26} & \raisebox{-1\normalbaselineskip}[0pt][0pt]{19} & \raisebox{-1\normalbaselineskip}[0pt][0pt]{B} & \raisebox{-1\normalbaselineskip}[0pt][0pt]{80} & \raisebox{-1\normalbaselineskip}[0pt][0pt]{19} & \raisebox{-1\normalbaselineskip}[0pt][0pt]{111} & \raisebox{-1\normalbaselineskip}[0pt][0pt]{2} & \raisebox{-1\normalbaselineskip}[0pt][0pt]{279} & \raisebox{-1\normalbaselineskip}[0pt][0pt]{4} & \raisebox{-1\normalbaselineskip}[0pt][0pt]{244} & \raisebox{-1\normalbaselineskip}[0pt][0pt]{24} & \raisebox{-1\normalbaselineskip}[0pt][0pt]{311} & \raisebox{-1\normalbaselineskip}[0pt][0pt]{\textit{4}} & \raisebox{-1\normalbaselineskip}[0pt][0pt]{3} \\
\cline{2-16}\end{tabular}\end{center}
{\small
\par\hangindent=12mm\hangafter=1 α Centauri pertinentes nescio quo modo translata sint infra, sub nr. 70. — In \textit{C} pro \textarabic{يط} 19′ cod. legendum \textarabic{لط} 39′; magnitudo cod. 1, quae cum falso nomine Pede Centauri convenit. \textit{P}, δ et \textit{h} inter se congruunt; \textit{A} cum catalogo ad annum 1211 reducto; \textit{B} et \textit{C} cum \textit{A}.
\par\hangindent=12mm\hangafter=1 17-24. Quoniam his circumpolaribus stellis \textit{P}, \textit{B}, \textit{C} desunt, impossibile fit numeros ex ipsorum inter se collatione emendare. Si cum catalogo fixarum comparantur, apparent hic illic gravissimae mendae, quarum naturam causamque investigare non possumus. Reliqui numeri, rudi quodam modo, inter se congruunt; nisi minuta essent corrupta, interdum \textit{A} anno 1211 conveniret. Itaque codicis numeros intactos servavi, inclinatis seu Italicis notis imprimens qui gravioribus scateant mendis. In nr. 22 minuta δ et \textit{h} aequalia esse deberent; discrepantia ab amanuense provenit (\textarabic{كر} 27, \textarabic{كو} 26).
\par\hangindent=12mm\hangafter=1 25. Multum δ et \textit{h} differunt ab eo quod praeberet \textit{P}, qui manifeste falsus est, et 116° 44′ circiter complecti debebat. Nam, hac quantitate, et data \textit{A} (cod. \textarabic{نا} 51′ pro \textarabic{يا} 11′) deducuntur numeri satis cum \textit{B} et \textit{C} congruentes; parvae tantum differentiae exstant quas non possumus explicare. In \textit{h} cod. \textarabic{قه} 105° pro \textarabic{فه} 85°.
\par\hangindent=12mm\hangafter=1 26. In \textit{C} cod. \textarabic{سمد} 344° pro \textarabic{سمه} 345°. Cum \textit{P} congruunt δ et \textit{h}; \textit{A} est 10′ minor quam quae a catalogo fixarum deducitur pro 1211, sed numero Albateniano exacte reperiuntur \textit{B} et \textit{C}.
\par\hangindent=12mm\hangafter=1 27. Magnitudo in cod. \textarabic{ب} 2 pro \textarabic{ـجـ} 3. \textit{A} cum catalogo ad annum 1211 reducto convenit; tunc proxime deducuntur \textit{B} (cuius vera quantitas 244° 29′) et \textit{C} (ubi 4′ error pro 35′).
\par}
\clearpage
\noindent\makebox[\textwidth]{\textbf{}\hfill{\small AL-BATTANI OPUS ASTRONOMICUM}\hfill\textbf{181}}\par\vspace{-1mm}\noindent\rule{\textwidth}{.4pt}\par\vspace{2mm}
\begin{center}\footnotesize\begin{tabular}{r|p{30mm}|r@{\hspace{1.5mm}}r|c|r@{\hspace{1.5mm}}r|r@{\hspace{1.5mm}}r|r@{\hspace{1.5mm}}r|r@{\hspace{1.5mm}}r|r@{\hspace{1.5mm}}r||c|}\cline{2-16}
 & \centering Stellarum nomina. & \multicolumn{2}{c|}{\rotatebox{90}{\parbox{14mm}{\centering\scriptsize Declinatio.}}} & \rotatebox{90}{\parbox{14mm}{\centering\scriptsize Plaga Caeli.}} & \multicolumn{2}{c|}{\rotatebox{90}{\parbox{24mm}{\centering\scriptsize Altitudo meridiana.}}} & \multicolumn{2}{c|}{\rotatebox{90}{\parbox{24mm}{\centering\scriptsize Dimidium commorationis supra Terram.}}} & \multicolumn{2}{c|}{\rotatebox{90}{\parbox{24mm}{\centering\scriptsize Gradus cum quibus culminant.}}} & \multicolumn{2}{c|}{\rotatebox{90}{\parbox{24mm}{\centering\scriptsize Gradus cum quibus oriuntur.}}} & \multicolumn{2}{c|}{\rotatebox{90}{\parbox{24mm}{\centering\scriptsize Gradus cum quibus occidunt.}}} & \rotatebox{90}{\parbox{14mm}{\centering\scriptsize Ordines magnitudinis.}} \\ \hline
28 & Latus Persei (α Per-\newline\hspace*{3mm}sei). & \raisebox{-1\normalbaselineskip}[0pt][0pt]{44°} & \raisebox{-1\normalbaselineskip}[0pt][0pt]{9′} & \raisebox{-1\normalbaselineskip}[0pt][0pt]{B} & \raisebox{-1\normalbaselineskip}[0pt][0pt]{98°} & \raisebox{-1\normalbaselineskip}[0pt][0pt]{9′} & \raisebox{-1\normalbaselineskip}[0pt][0pt]{134°} & \raisebox{-1\normalbaselineskip}[0pt][0pt]{52′} & \raisebox{-1\normalbaselineskip}[0pt][0pt]{34°} & \raisebox{-1\normalbaselineskip}[0pt][0pt]{30′} & \raisebox{-1\normalbaselineskip}[0pt][0pt]{340°} & \raisebox{-1\normalbaselineskip}[0pt][0pt]{0′} & \raisebox{-1\normalbaselineskip}[0pt][0pt]{63°} & \raisebox{-1\normalbaselineskip}[0pt][0pt]{0′} & \raisebox{-1\normalbaselineskip}[0pt][0pt]{2} \\
29 & Caput daemonis (Me-\newline\hspace*{3mm}dusae) in manu\newline\hspace*{3mm}Persei (β Persei). & \raisebox{-2\normalbaselineskip}[0pt][0pt]{36} & \raisebox{-2\normalbaselineskip}[0pt][0pt]{2} & \raisebox{-2\normalbaselineskip}[0pt][0pt]{B} & \raisebox{-2\normalbaselineskip}[0pt][0pt]{90} & \raisebox{-2\normalbaselineskip}[0pt][0pt]{2} & \raisebox{-2\normalbaselineskip}[0pt][0pt]{121} & \raisebox{-2\normalbaselineskip}[0pt][0pt]{55} & \raisebox{-2\normalbaselineskip}[0pt][0pt]{32} & \raisebox{-2\normalbaselineskip}[0pt][0pt]{12} & \raisebox{-2\normalbaselineskip}[0pt][0pt]{356} & \raisebox{-2\normalbaselineskip}[0pt][0pt]{53} & \raisebox{-2\normalbaselineskip}[0pt][0pt]{50} & \raisebox{-2\normalbaselineskip}[0pt][0pt]{43} & \raisebox{-2\normalbaselineskip}[0pt][0pt]{2} \\
30 & Capulus ensis (χ Per-\newline\hspace*{3mm}sei). & \raisebox{-1\normalbaselineskip}[0pt][0pt]{50} & \raisebox{-1\normalbaselineskip}[0pt][0pt]{39} & \raisebox{-1\normalbaselineskip}[0pt][0pt]{B} & \raisebox{-1\normalbaselineskip}[0pt][0pt]{104} & \raisebox{-1\normalbaselineskip}[0pt][0pt]{39} & \raisebox{-1\normalbaselineskip}[0pt][0pt]{152} & \raisebox{-1\normalbaselineskip}[0pt][0pt]{26} & \raisebox{-1\normalbaselineskip}[0pt][0pt]{17} & \raisebox{-1\normalbaselineskip}[0pt][0pt]{40} & \raisebox{-1\normalbaselineskip}[0pt][0pt]{295} & \raisebox{-1\normalbaselineskip}[0pt][0pt]{30} & \raisebox{-1\normalbaselineskip}[0pt][0pt]{64} & \raisebox{-1\normalbaselineskip}[0pt][0pt]{20} & \raisebox{-1\normalbaselineskip}[0pt][0pt]{neb.} \\
31 & Humerus Aurigae (β\newline\hspace*{3mm}Aurigae). & \raisebox{-1\normalbaselineskip}[0pt][0pt]{42} & \raisebox{-1\normalbaselineskip}[0pt][0pt]{17} & \raisebox{-1\normalbaselineskip}[0pt][0pt]{B} & \raisebox{-1\normalbaselineskip}[0pt][0pt]{96} & \raisebox{-1\normalbaselineskip}[0pt][0pt]{17} & \raisebox{-1\normalbaselineskip}[0pt][0pt]{131} & \raisebox{-1\normalbaselineskip}[0pt][0pt]{20} & \raisebox{-1\normalbaselineskip}[0pt][0pt]{69} & \raisebox{-1\normalbaselineskip}[0pt][0pt]{46} & \raisebox{-1\normalbaselineskip}[0pt][0pt]{40} & \raisebox{-1\normalbaselineskip}[0pt][0pt]{28} & \raisebox{-1\normalbaselineskip}[0pt][0pt]{90} & \raisebox{-1\normalbaselineskip}[0pt][0pt]{53} & \raisebox{-1\normalbaselineskip}[0pt][0pt]{2} \\
32 & \Name{an–Nasr aṭ–ṭā’ir}\newline\hspace*{3mm}« aquila volans »\newline\hspace*{3mm}(α Aquilae). & \raisebox{-2\normalbaselineskip}[0pt][0pt]{6} & \raisebox{-2\normalbaselineskip}[0pt][0pt]{22} & \raisebox{-2\normalbaselineskip}[0pt][0pt]{B} & \raisebox{-2\normalbaselineskip}[0pt][0pt]{60} & \raisebox{-2\normalbaselineskip}[0pt][0pt]{22} & \raisebox{-2\normalbaselineskip}[0pt][0pt]{94} & \raisebox{-2\normalbaselineskip}[0pt][0pt]{44} & \raisebox{-2\normalbaselineskip}[0pt][0pt]{282} & \raisebox{-2\normalbaselineskip}[0pt][0pt]{18} & \raisebox{-2\normalbaselineskip}[0pt][0pt]{261} & \raisebox{-2\normalbaselineskip}[0pt][0pt]{12} & \raisebox{-2\normalbaselineskip}[0pt][0pt]{301} & \raisebox{-2\normalbaselineskip}[0pt][0pt]{10} & \raisebox{-2\normalbaselineskip}[0pt][0pt]{2} \\
33 & \textit{Caput} Andromedae\newline\hspace*{3mm}(ο Andromedae). & \raisebox{-1\normalbaselineskip}[0pt][0pt]{\textit{37}} & \raisebox{-1\normalbaselineskip}[0pt][0pt]{\textit{14}} & \raisebox{-1\normalbaselineskip}[0pt][0pt]{B} & \raisebox{-1\normalbaselineskip}[0pt][0pt]{91} & \raisebox{-1\normalbaselineskip}[0pt][0pt]{14} & \raisebox{-1\normalbaselineskip}[0pt][0pt]{122} & \raisebox{-1\normalbaselineskip}[0pt][0pt]{27} & \raisebox{-1\normalbaselineskip}[0pt][0pt]{331} & \raisebox{-1\normalbaselineskip}[0pt][0pt]{22} & \raisebox{-1\normalbaselineskip}[0pt][0pt]{281} & \raisebox{-1\normalbaselineskip}[0pt][0pt]{54} & \raisebox{-1\normalbaselineskip}[0pt][0pt]{4} & \raisebox{-1\normalbaselineskip}[0pt][0pt]{52} & \raisebox{-1\normalbaselineskip}[0pt][0pt]{3} \\
34 & Ala Pegasi (α Pe-\newline\hspace*{3mm}gasi). & \raisebox{-1\normalbaselineskip}[0pt][0pt]{10} & \raisebox{-1\normalbaselineskip}[0pt][0pt]{16} & \raisebox{-1\normalbaselineskip}[0pt][0pt]{B} & \raisebox{-1\normalbaselineskip}[0pt][0pt]{64} & \raisebox{-1\normalbaselineskip}[0pt][0pt]{16} & \raisebox{-1\normalbaselineskip}[0pt][0pt]{97} & \raisebox{-1\normalbaselineskip}[0pt][0pt]{35} & \raisebox{-1\normalbaselineskip}[0pt][0pt]{330} & \raisebox{-1\normalbaselineskip}[0pt][0pt]{45} & \raisebox{-1\normalbaselineskip}[0pt][0pt]{309} & \raisebox{-1\normalbaselineskip}[0pt][0pt]{0} & \raisebox{-1\normalbaselineskip}[0pt][0pt]{343} & \raisebox{-1\normalbaselineskip}[0pt][0pt]{48} & \raisebox{-1\normalbaselineskip}[0pt][0pt]{2} \\
35 & Humerus Pegasi (β\newline\hspace*{3mm}Pegasi). & \raisebox{-1\normalbaselineskip}[0pt][0pt]{22} & \raisebox{-1\normalbaselineskip}[0pt][0pt]{31} & \raisebox{-1\normalbaselineskip}[0pt][0pt]{B} & \raisebox{-1\normalbaselineskip}[0pt][0pt]{76} & \raisebox{-1\normalbaselineskip}[0pt][0pt]{31} & \raisebox{-1\normalbaselineskip}[0pt][0pt]{107} & \raisebox{-1\normalbaselineskip}[0pt][0pt]{32} & \raisebox{-1\normalbaselineskip}[0pt][0pt]{330} & \raisebox{-1\normalbaselineskip}[0pt][0pt]{24} & \raisebox{-1\normalbaselineskip}[0pt][0pt]{296} & \raisebox{-1\normalbaselineskip}[0pt][0pt]{46} & \raisebox{-1\normalbaselineskip}[0pt][0pt]{351} & \raisebox{-1\normalbaselineskip}[0pt][0pt]{44} & \raisebox{-1\normalbaselineskip}[0pt][0pt]{2} \\
36 & Dorsus Pegasi (γ Pe-\newline\hspace*{3mm}gasi). & \raisebox{-1\normalbaselineskip}[0pt][0pt]{9} & \raisebox{-1\normalbaselineskip}[0pt][0pt]{7} & \raisebox{-1\normalbaselineskip}[0pt][0pt]{B} & \raisebox{-1\normalbaselineskip}[0pt][0pt]{63} & \raisebox{-1\normalbaselineskip}[0pt][0pt]{7} & \raisebox{-1\normalbaselineskip}[0pt][0pt]{96} & \raisebox{-1\normalbaselineskip}[0pt][0pt]{42} & \raisebox{-1\normalbaselineskip}[0pt][0pt]{348} & \raisebox{-1\normalbaselineskip}[0pt][0pt]{18} & \raisebox{-1\normalbaselineskip}[0pt][0pt]{332} & \raisebox{-1\normalbaselineskip}[0pt][0pt]{48} & \raisebox{-1\normalbaselineskip}[0pt][0pt]{356} & \raisebox{-1\normalbaselineskip}[0pt][0pt]{43} & \raisebox{-1\normalbaselineskip}[0pt][0pt]{2} \\
37 & Initium colli Serpen-\newline\hspace*{3mm}tis (β Serpentis). & \raisebox{-1\normalbaselineskip}[0pt][0pt]{20} & \raisebox{-1\normalbaselineskip}[0pt][0pt]{5} & \raisebox{-1\normalbaselineskip}[0pt][0pt]{B} & \raisebox{-1\normalbaselineskip}[0pt][0pt]{74} & \raisebox{-1\normalbaselineskip}[0pt][0pt]{5} & \raisebox{-1\normalbaselineskip}[0pt][0pt]{105} & \raisebox{-1\normalbaselineskip}[0pt][0pt]{24} & \raisebox{-1\normalbaselineskip}[0pt][0pt]{225} & \raisebox{-1\normalbaselineskip}[0pt][0pt]{27} & \raisebox{-1\normalbaselineskip}[0pt][0pt]{202} & \raisebox{-1\normalbaselineskip}[0pt][0pt]{42} & \raisebox{-1\normalbaselineskip}[0pt][0pt]{257} & \raisebox{-1\normalbaselineskip}[0pt][0pt]{20} & \raisebox{-1\normalbaselineskip}[0pt][0pt]{3} \\
\cline{2-16}\end{tabular}\end{center}
{\small
\par\hangindent=12mm\hangafter=1 28. In \textit{C} amanuensis maghrebinus in exemplari scriptura orientali seu nasḫī exarato invenerat \textarabic{شج} 303° pro \textarabic{سج} 63°; posuit igitur, scriptura maghrebina, \textarabic{سج} 303° loco \textarabic{صج} 63°. Cum \textit{P} congruunt δ et \textit{h}; \textit{A} convenit potius anno 1223 quam 1211, sed ab ea deducuntur bene \textit{B} et \textit{C}.
\par\hangindent=12mm\hangafter=1 29. In \textit{B} cod. \textarabic{سنج ج} 353° 3′ pro \textarabic{سنو نج} 356° 53′. \textit{P}, δ, \textit{h} bene inter se congruentes. \textit{A} (ubi minuta 12′ pro 6′) anno 1216 conveniret; sed ei \textit{B} et \textit{C} bene respondent.
\par\hangindent=12mm\hangafter=1 30. In \textit{P} cod. \textarabic{قيب} 112° pro \textarabic{قنب} 152°; in \textit{B} \textarabic{رصه} 265° pro \textarabic{رضه} 295°; in magnitudine 2 pro « nebulosa ». Proxime δ et \textit{h} cum \textit{P} congruunt; \textit{A} (17° 40′ pro 17° 17′) conveniret anno 1237, sed ei \textit{B} et \textit{C} bene respondent.
\par\hangindent=12mm\hangafter=1 31. \textit{P}, \textit{h}, δ bene congruentes; \textit{A} excedit 1° 38′ quantitatem ad annum 1211 pertinentem, sed exacte gignit \textit{B} et \textit{C} (ubi 53′ pro 51′ apparent).
\par\hangindent=12mm\hangafter=1 32. δ et \textit{h} differunt + 6′ ab iis quas \textit{P} reposcit, et − 4′ ab iis quae catalogo fixarum deducuntur pro 1211. \textit{A} e contra cum 1211 et \textit{B}, \textit{C}, optime congruit.
\par\hangindent=12mm\hangafter=1 33. Nomen falsum, de quo cfr. p. 154, nr. 3. Magnitudo in cod. \textarabic{ب} 2 pro \textarabic{ـجـ} 3. In \textit{h} cod. \textarabic{ضد} 94° pro \textarabic{ضا} 91°. Cum \textit{P} conveniret δ 36° 26′. Optime \textit{A} pro anno 1211, et \textit{B}, \textit{C}.
\par\hangindent=12mm\hangafter=1 34. In \textit{B} cod. \textarabic{صط} 69° pro \textarabic{سط} 309°; in \textit{C} cod. \textarabic{م} 40′ pro \textarabic{مح} 48′. A \textit{P} differt δ 5′ minutis. \textit{A} veritatem 24′ excedit, et anno 1240 respondet; sed a \textit{P} et \textit{A} codicis proxime \textit{B} et \textit{C} deducuntur.
\par\hangindent=12mm\hangafter=1 35. In \textit{C} cod. \textarabic{سند د} 354° 4′ pro \textarabic{سنا مد} 351° 44′. Reliqua optime.
\par\hangindent=12mm\hangafter=1 36. In \textit{A} cod. \textarabic{لح} 38′ pro \textarabic{ىح} = \textarabic{يح} 18′. A \textit{P} discrepat δ 3′. \textit{A} convenit anno 1218, sed cum \textit{B} et \textit{C} congruit.
\par\hangindent=12mm\hangafter=1 37. In \textit{P} cod. \textarabic{مو} 46′ pro \textarabic{كد} 24′, quae \textit{h} et δ reposcunt. Magnitudo cod. \textarabic{ب} 2 pro \textarabic{ـجـ} 3. \textit{A} congruit cum anno 1211; \textit{B} et \textit{C} cum \textit{A} et \textit{P} correcto.
\par}
\clearpage
\noindent\makebox[\textwidth]{\textbf{182}\hfill{\small C. A. NALLINO}\hfill\textbf{}}\par\vspace{-1mm}\noindent\rule{\textwidth}{.4pt}\par\vspace{2mm}
\begin{center}\footnotesize\begin{tabular}{r|p{30mm}|r@{\hspace{1.5mm}}r|c|r@{\hspace{1.5mm}}r|r@{\hspace{1.5mm}}r|r@{\hspace{1.5mm}}r|r@{\hspace{1.5mm}}r|r@{\hspace{1.5mm}}r||c|}\cline{2-16}
 & \centering Stellarum nomina. & \multicolumn{2}{c|}{\rotatebox{90}{\parbox{14mm}{\centering\scriptsize Declinatio.}}} & \rotatebox{90}{\parbox{14mm}{\centering\scriptsize Plaga Caeli.}} & \multicolumn{2}{c|}{\rotatebox{90}{\parbox{24mm}{\centering\scriptsize Altitudo meridiana.}}} & \multicolumn{2}{c|}{\rotatebox{90}{\parbox{24mm}{\centering\scriptsize Dimidium commorationis supra Terram.}}} & \multicolumn{2}{c|}{\rotatebox{90}{\parbox{24mm}{\centering\scriptsize Gradus cum quibus culminant.}}} & \multicolumn{2}{c|}{\rotatebox{90}{\parbox{24mm}{\centering\scriptsize Gradus cum quibus oriuntur.}}} & \multicolumn{2}{c|}{\rotatebox{90}{\parbox{24mm}{\centering\scriptsize Gradus cum quibus occidunt.}}} & \rotatebox{90}{\parbox{14mm}{\centering\scriptsize Ordines magnitudinis.}} \\ \hline
38 & Caput \Name{Afollon} («’Α-\newline\hspace*{3mm}πόλλων» i. e. Ge-\newline\hspace*{3mm}mini anterioris, α\newline\hspace*{3mm}Gemin.). & \raisebox{-3\normalbaselineskip}[0pt][0pt]{32°} & \raisebox{-3\normalbaselineskip}[0pt][0pt]{49′} & \raisebox{-3\normalbaselineskip}[0pt][0pt]{B} & \raisebox{-3\normalbaselineskip}[0pt][0pt]{106°} & \raisebox{-3\normalbaselineskip}[0pt][0pt]{49′} & \raisebox{-3\normalbaselineskip}[0pt][0pt]{117°} & \raisebox{-3\normalbaselineskip}[0pt][0pt]{57′} & \raisebox{-3\normalbaselineskip}[0pt][0pt]{95°} & \raisebox{-3\normalbaselineskip}[0pt][0pt]{10′} & \raisebox{-3\normalbaselineskip}[0pt][0pt]{86°} & \raisebox{-3\normalbaselineskip}[0pt][0pt]{28′} & \raisebox{-3\normalbaselineskip}[0pt][0pt]{104°} & \raisebox{-3\normalbaselineskip}[0pt][0pt]{32′} & \raisebox{-3\normalbaselineskip}[0pt][0pt]{2} \\
39 & Caput \Name{Īraqlās} («‘Η-\newline\hspace*{3mm}ρακλῆς» i. e. Ge-\newline\hspace*{3mm}mini posterioris, β\newline\hspace*{3mm}Gemin.). & \raisebox{-3\normalbaselineskip}[0pt][0pt]{29} & \raisebox{-3\normalbaselineskip}[0pt][0pt]{32} & \raisebox{-3\normalbaselineskip}[0pt][0pt]{B} & \raisebox{-3\normalbaselineskip}[0pt][0pt]{103} & \raisebox{-3\normalbaselineskip}[0pt][0pt]{32} & \raisebox{-3\normalbaselineskip}[0pt][0pt]{114} & \raisebox{-3\normalbaselineskip}[0pt][0pt]{18} & \raisebox{-3\normalbaselineskip}[0pt][0pt]{98} & \raisebox{-3\normalbaselineskip}[0pt][0pt]{33} & \raisebox{-3\normalbaselineskip}[0pt][0pt]{93} & \raisebox{-3\normalbaselineskip}[0pt][0pt]{12} & \raisebox{-3\normalbaselineskip}[0pt][0pt]{104} & \raisebox{-3\normalbaselineskip}[0pt][0pt]{35} & \raisebox{-3\normalbaselineskip}[0pt][0pt]{2} \\
40 & Medium Pleiadum,\newline\hspace*{3mm}nebulosa (η Tauri). & \raisebox{-1\normalbaselineskip}[0pt][0pt]{19} & \raisebox{-1\normalbaselineskip}[0pt][0pt]{18} & \raisebox{-1\normalbaselineskip}[0pt][0pt]{B} & \raisebox{-1\normalbaselineskip}[0pt][0pt]{73} & \raisebox{-1\normalbaselineskip}[0pt][0pt]{18} & \raisebox{-1\normalbaselineskip}[0pt][0pt]{104} & \raisebox{-1\normalbaselineskip}[0pt][0pt]{45} & \raisebox{-1\normalbaselineskip}[0pt][0pt]{43} & \raisebox{-1\normalbaselineskip}[0pt][0pt]{13} & \raisebox{-1\normalbaselineskip}[0pt][0pt]{39} & \raisebox{-1\normalbaselineskip}[0pt][0pt]{22} & \raisebox{-1\normalbaselineskip}[0pt][0pt]{44} & \raisebox{-1\normalbaselineskip}[0pt][0pt]{17} & \raisebox{-1\normalbaselineskip}[0pt][0pt]{4} \\
41 & Pectus Cancri, ne-\newline\hspace*{3mm}bulosa (Praesepe). & \raisebox{-1\normalbaselineskip}[0pt][0pt]{22} & \raisebox{-1\normalbaselineskip}[0pt][0pt]{10} & \raisebox{-1\normalbaselineskip}[0pt][0pt]{B} & \raisebox{-1\normalbaselineskip}[0pt][0pt]{76} & \raisebox{-1\normalbaselineskip}[0pt][0pt]{10} & \raisebox{-1\normalbaselineskip}[0pt][0pt]{107} & \raisebox{-1\normalbaselineskip}[0pt][0pt]{15} & \raisebox{-1\normalbaselineskip}[0pt][0pt]{111} & \raisebox{-1\normalbaselineskip}[0pt][0pt]{54} & \raisebox{-1\normalbaselineskip}[0pt][0pt]{111} & \raisebox{-1\normalbaselineskip}[0pt][0pt]{36} & \raisebox{-1\normalbaselineskip}[0pt][0pt]{112} & \raisebox{-1\normalbaselineskip}[0pt][0pt]{16} & \raisebox{-1\normalbaselineskip}[0pt][0pt]{neb.} \\
42 & Dorsus Leonis (δ Leo-\newline\hspace*{3mm}nis). & \raisebox{-1\normalbaselineskip}[0pt][0pt]{25} & \raisebox{-1\normalbaselineskip}[0pt][0pt]{2} & \raisebox{-1\normalbaselineskip}[0pt][0pt]{B} & \raisebox{-1\normalbaselineskip}[0pt][0pt]{79} & \raisebox{-1\normalbaselineskip}[0pt][0pt]{2} & \raisebox{-1\normalbaselineskip}[0pt][0pt]{109} & \raisebox{-1\normalbaselineskip}[0pt][0pt]{51} & \raisebox{-1\normalbaselineskip}[0pt][0pt]{151} & \raisebox{-1\normalbaselineskip}[0pt][0pt]{0} & \raisebox{-1\normalbaselineskip}[0pt][0pt]{141} & \raisebox{-1\normalbaselineskip}[0pt][0pt]{34} & \raisebox{-1\normalbaselineskip}[0pt][0pt]{168} & \raisebox{-1\normalbaselineskip}[0pt][0pt]{44} & \raisebox{-1\normalbaselineskip}[0pt][0pt]{2} \\
43 & \textit{Caput Orionis} (α\newline\hspace*{3mm}Andromedae). & \raisebox{-1\normalbaselineskip}[0pt][0pt]{19} & \raisebox{-1\normalbaselineskip}[0pt][0pt]{38} & \raisebox{-1\normalbaselineskip}[0pt][0pt]{B} & \raisebox{-1\normalbaselineskip}[0pt][0pt]{73} & \raisebox{-1\normalbaselineskip}[0pt][0pt]{38} & \raisebox{-1\normalbaselineskip}[0pt][0pt]{105} & \raisebox{-1\normalbaselineskip}[0pt][0pt]{2} & \raisebox{-1\normalbaselineskip}[0pt][0pt]{347} & \raisebox{-1\normalbaselineskip}[0pt][0pt]{15} & \raisebox{-1\normalbaselineskip}[0pt][0pt]{319} & \raisebox{-1\normalbaselineskip}[0pt][0pt]{30} & \raisebox{-1\normalbaselineskip}[0pt][0pt]{2} & \raisebox{-1\normalbaselineskip}[0pt][0pt]{44} & \raisebox{-1\normalbaselineskip}[0pt][0pt]{3} \\
44 & \textit{Mulier (Cassiopea)}\newline\hspace*{3mm}\textit{quae vocatur al-}\newline\hspace*{3mm}\textit{–kaff al–khaḍīb}\newline\hspace*{3mm}\textit{« manus tincta »}\newline\hspace*{3mm}(α Ophiuchi). & \raisebox{-4\normalbaselineskip}[0pt][0pt]{14} & \raisebox{-4\normalbaselineskip}[0pt][0pt]{18} & \raisebox{-4\normalbaselineskip}[0pt][0pt]{B} & \raisebox{-4\normalbaselineskip}[0pt][0pt]{68} & \raisebox{-4\normalbaselineskip}[0pt][0pt]{18} & \raisebox{-4\normalbaselineskip}[0pt][0pt]{100} & \raisebox{-4\normalbaselineskip}[0pt][0pt]{38} & \raisebox{-4\normalbaselineskip}[0pt][0pt]{251} & \raisebox{-4\normalbaselineskip}[0pt][0pt]{56} & \raisebox{-4\normalbaselineskip}[0pt][0pt]{228} & \raisebox{-4\normalbaselineskip}[0pt][0pt]{58} & \raisebox{-4\normalbaselineskip}[0pt][0pt]{278} & \raisebox{-4\normalbaselineskip}[0pt][0pt]{34} & \raisebox{-4\normalbaselineskip}[0pt][0pt]{3} \\
\cline{2-16}\end{tabular}\end{center}
{\small
\par\hangindent=12mm\hangafter=1 38. Nomen \textit{Afollon} in catalogo fixarum non memoratur. In \textit{P} cod. \textarabic{يز} 17′ pro \textarabic{نز} 57′; in \textit{B} cod. \textarabic{قو} 106° pro \textarabic{فو} 86°. Omnia optime; \textit{A} anno 1211 convenit.
\par\hangindent=12mm\hangafter=1 39. Nomen \textit{Īraqlās} in catalogo non memoratur. In \textit{C} cod. \textarabic{ل} 30′ pro \textarabic{له} 35′. Omnia optime pro anno 1211.
\par\hangindent=12mm\hangafter=1 40. In δ et \textit{h} cod. \textarabic{لح} 38′ pro \textarabic{ىح} = \textarabic{يح} 18′; in \textit{B} cod. \textarabic{نط} 59′ pro \textarabic{لط} 39′; in \textit{C} cod. \textarabic{قمد ل} 144° 30′ pro \textarabic{مد ىر} 44° 17′; magnitudo cod. \textarabic{ـجـ} 3 pro \textarabic{د} 4. Pro anno 1211 debebant esse δ 20° 18′,8, \textit{P} 105° 36′, \textit{A} 42° 10′,7; sed e datis numeris bene \textit{B} et \textit{C} proveniunt.
\par\hangindent=12mm\hangafter=1 41. In \textit{P} cod. \textarabic{مه} 45′ pro \textarabic{ىه} = \textarabic{يه} 15′; in magnitudine cod. 3 pro « nebulosa ». Omnia optime inter se congruentia; \textit{A} pro anno 1211. Sed iuxta catalogum fixarum δ 22° 27′,8 esse debebat.
\par\hangindent=12mm\hangafter=1 42. In \textit{A} cod. \textarabic{قيا} 111° pro \textarabic{قنا} 151°; in \textit{B} cod. 30′ pro 34′; in \textit{C} cod. 40′ pro 44′. Omnia optime; \textit{A} pro anno 1211. Sed iuxta catalogum δ 25° 50′,7 debebat esse.
\par\hangindent=12mm\hangafter=1 43. Nomen falsum; caput Orionis infra, nr. 54, occurrit. Stella cui elementa hic exposita satis bene conveniunt est Caput Andromedae, prima stellarum Pegasi in catalogo p. 153, et falso supra (nr. 33) indicata. Calculis, iuxta catalogi elementa, pro anno 1211 ductis, pauca tantum minuta differentiae in \textit{A}, sed 3° 35′ in δ colliguntur. At δ codicis confirmatur arcu semidiurno \textit{P}; a quo et ab \textit{A} codicis optime \textit{B} et \textit{C} proveniunt. Elementa igitur, inter se congruentia, pertinent ad locum Caeli circiter 3° $\tfrac{1}{2}$ ab austro α \textit{Andromedae}, in quo nulla stella priorum magnitudinum conspicitur. In \textit{C} cod. \textarabic{رب} 202° pro \textarabic{ب} 2°.
\par\hangindent=12mm\hangafter=1 44. Nomen falsum, nam est prima stellarum Ophiuchi p. 154. Quantitate + 3′ discrepant δ et \textit{h} a dato \textit{P} a quo et ab \textit{A} (anno 1211 convenienti) bene \textit{B} et \textit{C} proveniunt. In \textit{B} cod. 50′ pro 58′.
\par}
\clearpage
\noindent\makebox[\textwidth]{\textbf{}\hfill{\small AL-BATTANI OPUS ASTRONOMICUM}\hfill\textbf{183}}\par\vspace{-1mm}\noindent\rule{\textwidth}{.4pt}\par\vspace{2mm}
\begin{center}\footnotesize\begin{tabular}{r|p{30mm}|r@{\hspace{1.5mm}}r|c|r@{\hspace{1.5mm}}r|r@{\hspace{1.5mm}}r|r@{\hspace{1.5mm}}r|r@{\hspace{1.5mm}}r|r@{\hspace{1.5mm}}r||c|}\cline{2-16}
 & \centering Stellarum nomina. & \multicolumn{2}{c|}{\rotatebox{90}{\parbox{14mm}{\centering\scriptsize Declinatio.}}} & \rotatebox{90}{\parbox{14mm}{\centering\scriptsize Plaga Caeli.}} & \multicolumn{2}{c|}{\rotatebox{90}{\parbox{24mm}{\centering\scriptsize Altitudo meridiana.}}} & \multicolumn{2}{c|}{\rotatebox{90}{\parbox{24mm}{\centering\scriptsize Dimidium commorationis supra Terram.}}} & \multicolumn{2}{c|}{\rotatebox{90}{\parbox{24mm}{\centering\scriptsize Gradus cum quibus culminant.}}} & \multicolumn{2}{c|}{\rotatebox{90}{\parbox{24mm}{\centering\scriptsize Gradus cum quibus oriuntur.}}} & \multicolumn{2}{c|}{\rotatebox{90}{\parbox{24mm}{\centering\scriptsize Gradus cum quibus occidunt.}}} & \rotatebox{90}{\parbox{14mm}{\centering\scriptsize Ordines magnitudinis.}} \\ \hline
45 & Lanx australis Li-\newline\hspace*{3mm}brae (α Librae). & \raisebox{-1\normalbaselineskip}[0pt][0pt]{8°} & \raisebox{-1\normalbaselineskip}[0pt][0pt]{42′} & \raisebox{-1\normalbaselineskip}[0pt][0pt]{A} & \raisebox{-1\normalbaselineskip}[0pt][0pt]{45°} & \raisebox{-1\normalbaselineskip}[0pt][0pt]{18′} & \raisebox{-1\normalbaselineskip}[0pt][0pt]{83°} & \raisebox{-1\normalbaselineskip}[0pt][0pt]{38′} & \raisebox{-1\normalbaselineskip}[0pt][0pt]{210°} & \raisebox{-1\normalbaselineskip}[0pt][0pt]{14′} & \raisebox{-1\normalbaselineskip}[0pt][0pt]{208°} & \raisebox{-1\normalbaselineskip}[0pt][0pt]{25′} & \raisebox{-1\normalbaselineskip}[0pt][0pt]{213°} & \raisebox{-1\normalbaselineskip}[0pt][0pt]{27′} & \raisebox{-1\normalbaselineskip}[0pt][0pt]{2} \\
46 & Lanx borealis Librae\newline\hspace*{3mm}(β Librae). & \raisebox{-1\normalbaselineskip}[0pt][0pt]{3} & \raisebox{-1\normalbaselineskip}[0pt][0pt]{43} & \raisebox{-1\normalbaselineskip}[0pt][0pt]{A} & \raisebox{-1\normalbaselineskip}[0pt][0pt]{50} & \raisebox{-1\normalbaselineskip}[0pt][0pt]{17} & \raisebox{-1\normalbaselineskip}[0pt][0pt]{87} & \raisebox{-1\normalbaselineskip}[0pt][0pt]{18} & \raisebox{-1\normalbaselineskip}[0pt][0pt]{216} & \raisebox{-1\normalbaselineskip}[0pt][0pt]{35} & \raisebox{-1\normalbaselineskip}[0pt][0pt]{210} & \raisebox{-1\normalbaselineskip}[0pt][0pt]{32} & \raisebox{-1\normalbaselineskip}[0pt][0pt]{227} & \raisebox{-1\normalbaselineskip}[0pt][0pt]{0} & \raisebox{-1\normalbaselineskip}[0pt][0pt]{2} \\
47 & Cor Scorpionis (α\newline\hspace*{3mm}Scorpii). & \raisebox{-1\normalbaselineskip}[0pt][0pt]{22} & \raisebox{-1\normalbaselineskip}[0pt][0pt]{10} & \raisebox{-1\normalbaselineskip}[0pt][0pt]{A} & \raisebox{-1\normalbaselineskip}[0pt][0pt]{31} & \raisebox{-1\normalbaselineskip}[0pt][0pt]{50} & \raisebox{-1\normalbaselineskip}[0pt][0pt]{72} & \raisebox{-1\normalbaselineskip}[0pt][0pt]{40} & \raisebox{-1\normalbaselineskip}[0pt][0pt]{229} & \raisebox{-1\normalbaselineskip}[0pt][0pt]{15} & \raisebox{-1\normalbaselineskip}[0pt][0pt]{232} & \raisebox{-1\normalbaselineskip}[0pt][0pt]{27} & \raisebox{-1\normalbaselineskip}[0pt][0pt]{224} & \raisebox{-1\normalbaselineskip}[0pt][0pt]{5} & \raisebox{-1\normalbaselineskip}[0pt][0pt]{2} \\
48 & Malleolus Sagittarii\newline\hspace*{3mm}(β Sagittarii). & \raisebox{-1\normalbaselineskip}[0pt][0pt]{46} & \raisebox{-1\normalbaselineskip}[0pt][0pt]{34} & \raisebox{-1\normalbaselineskip}[0pt][0pt]{A} & \raisebox{-1\normalbaselineskip}[0pt][0pt]{7} & \raisebox{-1\normalbaselineskip}[0pt][0pt]{26} & \raisebox{-1\normalbaselineskip}[0pt][0pt]{39} & \raisebox{-1\normalbaselineskip}[0pt][0pt]{52} & \raisebox{-1\normalbaselineskip}[0pt][0pt]{268} & \raisebox{-1\normalbaselineskip}[0pt][0pt]{28} & \raisebox{-1\normalbaselineskip}[0pt][0pt]{300} & \raisebox{-1\normalbaselineskip}[0pt][0pt]{52} & \raisebox{-1\normalbaselineskip}[0pt][0pt]{235} & \raisebox{-1\normalbaselineskip}[0pt][0pt]{18} & \raisebox{-1\normalbaselineskip}[0pt][0pt]{3} \\
49 & Quae sequitur acu-\newline\hspace*{3mm}leum Scorpionis,\newline\hspace*{3mm}nebulosa. & \raisebox{-2\normalbaselineskip}[0pt][0pt]{35} & \raisebox{-2\normalbaselineskip}[0pt][0pt]{23} & \raisebox{-2\normalbaselineskip}[0pt][0pt]{A} & \raisebox{-2\normalbaselineskip}[0pt][0pt]{18} & \raisebox{-2\normalbaselineskip}[0pt][0pt]{37} & \raisebox{-2\normalbaselineskip}[0pt][0pt]{58} & \raisebox{-2\normalbaselineskip}[0pt][0pt]{56} & \raisebox{-2\normalbaselineskip}[0pt][0pt]{250} & \raisebox{-2\normalbaselineskip}[0pt][0pt]{40} & \raisebox{-2\normalbaselineskip}[0pt][0pt]{262} & \raisebox{-2\normalbaselineskip}[0pt][0pt]{27} & \raisebox{-2\normalbaselineskip}[0pt][0pt]{235} & \raisebox{-2\normalbaselineskip}[0pt][0pt]{0} & \raisebox{-2\normalbaselineskip}[0pt][0pt]{neb.} \\
50 & Oculus Sagittarii, ne-\newline\hspace*{3mm}bulosa (ν² Sagitta-\newline\hspace*{3mm}rii). & \raisebox{-2\normalbaselineskip}[0pt][0pt]{22} & \raisebox{-2\normalbaselineskip}[0pt][0pt]{46} & \raisebox{-2\normalbaselineskip}[0pt][0pt]{A} & \raisebox{-2\normalbaselineskip}[0pt][0pt]{31} & \raisebox{-2\normalbaselineskip}[0pt][0pt]{14} & \raisebox{-2\normalbaselineskip}[0pt][0pt]{72} & \raisebox{-2\normalbaselineskip}[0pt][0pt]{16} & \raisebox{-2\normalbaselineskip}[0pt][0pt]{266} & \raisebox{-2\normalbaselineskip}[0pt][0pt]{39} & \raisebox{-2\normalbaselineskip}[0pt][0pt]{266} & \raisebox{-2\normalbaselineskip}[0pt][0pt]{0} & \raisebox{-2\normalbaselineskip}[0pt][0pt]{267} & \raisebox{-2\normalbaselineskip}[0pt][0pt]{20} & \raisebox{-2\normalbaselineskip}[0pt][0pt]{3} \\
51 & Cuspis sagittae Sagit-\newline\hspace*{3mm}tarii (γ Sagittarii). & \raisebox{-1\normalbaselineskip}[0pt][0pt]{29} & \raisebox{-1\normalbaselineskip}[0pt][0pt]{8} & \raisebox{-1\normalbaselineskip}[0pt][0pt]{A} & \raisebox{-1\normalbaselineskip}[0pt][0pt]{24} & \raisebox{-1\normalbaselineskip}[0pt][0pt]{52} & \raisebox{-1\normalbaselineskip}[0pt][0pt]{66} & \raisebox{-1\normalbaselineskip}[0pt][0pt]{8} & \raisebox{-1\normalbaselineskip}[0pt][0pt]{255} & \raisebox{-1\normalbaselineskip}[0pt][0pt]{17} & \raisebox{-1\normalbaselineskip}[0pt][0pt]{260} & \raisebox{-1\normalbaselineskip}[0pt][0pt]{31} & \raisebox{-1\normalbaselineskip}[0pt][0pt]{248} & \raisebox{-1\normalbaselineskip}[0pt][0pt]{52} & \raisebox{-1\normalbaselineskip}[0pt][0pt]{3} \\
52 & Coma Leonis, crines\newline\hspace*{3mm}plexi (Prima in Co-\newline\hspace*{3mm}ma Berenices). & \raisebox{-2\normalbaselineskip}[0pt][0pt]{36} & \raisebox{-2\normalbaselineskip}[0pt][0pt]{10} & \raisebox{-2\normalbaselineskip}[0pt][0pt]{B} & \raisebox{-2\normalbaselineskip}[0pt][0pt]{90} & \raisebox{-2\normalbaselineskip}[0pt][0pt]{10} & \raisebox{-2\normalbaselineskip}[0pt][0pt]{121} & \raisebox{-2\normalbaselineskip}[0pt][0pt]{58} & \raisebox{-2\normalbaselineskip}[0pt][0pt]{170} & \raisebox{-2\normalbaselineskip}[0pt][0pt]{43} & \raisebox{-2\normalbaselineskip}[0pt][0pt]{146} & \raisebox{-2\normalbaselineskip}[0pt][0pt]{40} & \raisebox{-2\normalbaselineskip}[0pt][0pt]{215} & \raisebox{-2\normalbaselineskip}[0pt][0pt]{50} & \raisebox{-2\normalbaselineskip}[0pt][0pt]{obsc.} \\
\cline{2-16}\end{tabular}\end{center}
{\small
\par\hangindent=12mm\hangafter=1 45. 45 (1). Omnia optime; \textit{A} pro 1211. In \textit{P} cod. \textarabic{قج} 103° pro \textarabic{فج} 83°; in \textit{B} \textarabic{له} 35′ pro \textarabic{كه} 25′.
\par\hangindent=12mm\hangafter=1 46. In δ cod. \textarabic{ح مز} 8° 47′ pro \textarabic{ج مج} 3° 43′; in \textit{h} \textarabic{لز} 37′ pro \textarabic{ىز} = \textarabic{يز} 17′; in \textit{P} \textarabic{قز} 107° pro \textarabic{فز} 87°. Omnia satis bene; \textit{A} pro 1211.
\par\hangindent=12mm\hangafter=1 47. In elementis huius stellae supputandis, ut in nr. 43, auctor situ falso usus est, longitudine perperam 230° 0′ pro 233° 50′ sumpta. Quantitate − 7′,8 discrepat δ ab eo quod \textit{P} praeberet; a \textit{P} et \textit{A} (anno 1211 satis convenienti) optime \textit{B} et \textit{C} deducuntur. In \textit{A} cod. \textarabic{نه} 55′ pro \textarabic{يه} 15′; in \textit{C} \textarabic{يه} 15′ pro \textarabic{ه} 5′.
\par\hangindent=12mm\hangafter=1 48. Una e paucis stellis in quibus omnia, etiam δ et \textit{h}, ad amussim inter se et cum catalogo ad 1211 reducto congruunt.
\par\hangindent=12mm\hangafter=1 49. In \textit{P} cod. \textarabic{يح} 18° pro \textarabic{نح} 58°; in \textit{h} \textarabic{يج} 13° pro \textarabic{يح} 18°. Omnia optime; \textit{A} pro 1211.
\par\hangindent=12mm\hangafter=1 50. In \textit{h} cod. \textarabic{ند} 54′ pro \textarabic{يد} 14′; in \textit{B} cod. \textarabic{رصز} 267° pro \textarabic{رصو} 266°. Omnia optime; \textit{A} pro 1211.
\par\hangindent=12mm\hangafter=1 51. In δ cod. \textarabic{\AbjadZero{}} 0′ pro \textarabic{ح} 8′; in \textit{h} \textarabic{\AbjadZero{} كه} 25° 0′ pro \textarabic{كد نب} 24° 52′; in \textit{A} \textarabic{مز} 47′ pro \textarabic{ىز} = \textarabic{يز} 17′. Omnia optime, etiam δ et \textit{h}, non solum inter se sed etiam cum catalogo pro 1211 reducto.
\par\hangindent=12mm\hangafter=1 52. In \textit{P} cod. \textarabic{يح} 18′ pro \textarabic{نح} 58′; in \textit{B} et \textit{C} minuta rotunda (40′ et 50′ pro 41′ et 53′); in \textit{C} cod. \textarabic{رنه} 255° pro \textarabic{ريه} 215°. Ab eo quod arcui semidiurno conveniret discrepant δ et \textit{h} 7′; \textit{A} est 8′ minor quam quae anno 1211 responderet. Datis \textit{A} et \textit{P} bene \textit{B} et \textit{C} deducuntur.
\par\hangindent=12mm\hangafter=1 53. Fortasse vox \textarabic{العقرب} « Scorpio » corrupta ex \textarabic{الضفيرة} « Crines plexi » i. e. Coma Berenices. In δ minuta sunt rotunda facta, 40′ nempe pro 44′; unde \textit{h} quoque 40′ pro 44′. In \textit{A} cod. \textarabic{ند} 54′ pro \textarabic{يد} 14′; quibus correctis, optime respondet \textit{B}, \textit{C} et anno 1211.
\par}
\par\vspace{2mm}\begin{center}\rule{40mm}{.4pt}\end{center}{\small (1) Huius stellae, a qua fol. 238,v. codicis incipit, et sequentium non indicatur in codice magnitudo; omnesque, iuxta titulum folii, tertiae magnitudinis dicuntur. Veras quantitates restitui ad fidem catalogi pro anno 1191.}
\clearpage
\noindent\makebox[\textwidth]{\textbf{184}\hfill{\small C. A. NALLINO}\hfill\textbf{}}\par\vspace{-1mm}\noindent\rule{\textwidth}{.4pt}\par\vspace{2mm}
\begin{center}\footnotesize\begin{tabular}{r|p{30mm}|r@{\hspace{1.5mm}}r|c|r@{\hspace{1.5mm}}r|r@{\hspace{1.5mm}}r|r@{\hspace{1.5mm}}r|r@{\hspace{1.5mm}}r|r@{\hspace{1.5mm}}r||c|}\cline{2-16}
 & \centering Stellarum nomina. & \multicolumn{2}{c|}{\rotatebox{90}{\parbox{14mm}{\centering\scriptsize Declinatio.}}} & \rotatebox{90}{\parbox{14mm}{\centering\scriptsize Plaga Caeli.}} & \multicolumn{2}{c|}{\rotatebox{90}{\parbox{24mm}{\centering\scriptsize Altitudo meridiana.}}} & \multicolumn{2}{c|}{\rotatebox{90}{\parbox{24mm}{\centering\scriptsize Dimidium commorationis supra Terram.}}} & \multicolumn{2}{c|}{\rotatebox{90}{\parbox{24mm}{\centering\scriptsize Gradus cum quibus culminant.}}} & \multicolumn{2}{c|}{\rotatebox{90}{\parbox{24mm}{\centering\scriptsize Gradus cum quibus oriuntur.}}} & \multicolumn{2}{c|}{\rotatebox{90}{\parbox{24mm}{\centering\scriptsize Gradus cum quibus occidunt.}}} & \rotatebox{90}{\parbox{14mm}{\centering\scriptsize Ordines magnitudinis.}} \\ \hline
53 & Una duarum obscu-\newline\hspace*{3mm}rarum quae sunt\newline\hspace*{3mm}in \textit{Scorpione} (Se-\newline\hspace*{3mm}cunda Comae Be-\newline\hspace*{3mm}renices). & \raisebox{-4\normalbaselineskip}[0pt][0pt]{31°} & \raisebox{-4\normalbaselineskip}[0pt][0pt]{40′} & \raisebox{-4\normalbaselineskip}[0pt][0pt]{B} & \raisebox{-4\normalbaselineskip}[0pt][0pt]{85°} & \raisebox{-4\normalbaselineskip}[0pt][0pt]{40′} & \raisebox{-4\normalbaselineskip}[0pt][0pt]{116°} & \raisebox{-4\normalbaselineskip}[0pt][0pt]{41′} & \raisebox{-4\normalbaselineskip}[0pt][0pt]{167°} & \raisebox{-4\normalbaselineskip}[0pt][0pt]{14′} & \raisebox{-4\normalbaselineskip}[0pt][0pt]{148°} & \raisebox{-4\normalbaselineskip}[0pt][0pt]{22′} & \raisebox{-4\normalbaselineskip}[0pt][0pt]{203°} & \raisebox{-4\normalbaselineskip}[0pt][0pt]{29′} & \raisebox{-4\normalbaselineskip}[0pt][0pt]{obsc.} \\
54 & Caput Orionis (λ O-\newline\hspace*{3mm}rionis). & \raisebox{-1\normalbaselineskip}[0pt][0pt]{5} & \raisebox{-1\normalbaselineskip}[0pt][0pt]{16} & \raisebox{-1\normalbaselineskip}[0pt][0pt]{B} & \raisebox{-1\normalbaselineskip}[0pt][0pt]{59} & \raisebox{-1\normalbaselineskip}[0pt][0pt]{16} & \raisebox{-1\normalbaselineskip}[0pt][0pt]{93} & \raisebox{-1\normalbaselineskip}[0pt][0pt]{46} & \raisebox{-1\normalbaselineskip}[0pt][0pt]{71} & \raisebox{-1\normalbaselineskip}[0pt][0pt]{0} & \raisebox{-1\normalbaselineskip}[0pt][0pt]{84} & \raisebox{-1\normalbaselineskip}[0pt][0pt]{30} & \raisebox{-1\normalbaselineskip}[0pt][0pt]{59} & \raisebox{-1\normalbaselineskip}[0pt][0pt]{45} & \raisebox{-1\normalbaselineskip}[0pt][0pt]{neb.} \\
55 & Humerus sinister O-\newline\hspace*{3mm}rionis (γ Orionis). & \raisebox{-1\normalbaselineskip}[0pt][0pt]{3} & \raisebox{-1\normalbaselineskip}[0pt][0pt]{47} & \raisebox{-1\normalbaselineskip}[0pt][0pt]{B} & \raisebox{-1\normalbaselineskip}[0pt][0pt]{57} & \raisebox{-1\normalbaselineskip}[0pt][0pt]{47} & \raisebox{-1\normalbaselineskip}[0pt][0pt]{92} & \raisebox{-1\normalbaselineskip}[0pt][0pt]{44} & \raisebox{-1\normalbaselineskip}[0pt][0pt]{68} & \raisebox{-1\normalbaselineskip}[0pt][0pt]{22} & \raisebox{-1\normalbaselineskip}[0pt][0pt]{82} & \raisebox{-1\normalbaselineskip}[0pt][0pt]{50} & \raisebox{-1\normalbaselineskip}[0pt][0pt]{56} & \raisebox{-1\normalbaselineskip}[0pt][0pt]{42} & \raisebox{-1\normalbaselineskip}[0pt][0pt]{2} \\
56 & Prima ex iis quae\newline\hspace*{3mm}sunt in eius zona\newline\hspace*{3mm}(δ Orionis). & \raisebox{-2\normalbaselineskip}[0pt][0pt]{2} & \raisebox{-2\normalbaselineskip}[0pt][0pt]{32} & \raisebox{-2\normalbaselineskip}[0pt][0pt]{A} & \raisebox{-2\normalbaselineskip}[0pt][0pt]{51} & \raisebox{-2\normalbaselineskip}[0pt][0pt]{28} & \raisebox{-2\normalbaselineskip}[0pt][0pt]{88} & \raisebox{-2\normalbaselineskip}[0pt][0pt]{11} & \raisebox{-2\normalbaselineskip}[0pt][0pt]{70} & \raisebox{-2\normalbaselineskip}[0pt][0pt]{33} & \raisebox{-2\normalbaselineskip}[0pt][0pt]{89} & \raisebox{-2\normalbaselineskip}[0pt][0pt]{18} & \raisebox{-2\normalbaselineskip}[0pt][0pt]{54} & \raisebox{-2\normalbaselineskip}[0pt][0pt]{54} & \raisebox{-2\normalbaselineskip}[0pt][0pt]{2} \\
57 & Media earum (ε O-\newline\hspace*{3mm}rionis). & \raisebox{-1\normalbaselineskip}[0pt][0pt]{2} & \raisebox{-1\normalbaselineskip}[0pt][0pt]{54} & \raisebox{-1\normalbaselineskip}[0pt][0pt]{A} & \raisebox{-1\normalbaselineskip}[0pt][0pt]{51} & \raisebox{-1\normalbaselineskip}[0pt][0pt]{6} & \raisebox{-1\normalbaselineskip}[0pt][0pt]{87} & \raisebox{-1\normalbaselineskip}[0pt][0pt]{55} & \raisebox{-1\normalbaselineskip}[0pt][0pt]{72} & \raisebox{-1\normalbaselineskip}[0pt][0pt]{20} & \raisebox{-1\normalbaselineskip}[0pt][0pt]{91} & \raisebox{-1\normalbaselineskip}[0pt][0pt]{40} & \raisebox{-1\normalbaselineskip}[0pt][0pt]{55} & \raisebox{-1\normalbaselineskip}[0pt][0pt]{53} & \raisebox{-1\normalbaselineskip}[0pt][0pt]{2} \\
58 & Tertia earum (ζ O-\newline\hspace*{3mm}rionis). & \raisebox{-1\normalbaselineskip}[0pt][0pt]{3} & \raisebox{-1\normalbaselineskip}[0pt][0pt]{35} & \raisebox{-1\normalbaselineskip}[0pt][0pt]{A} & \raisebox{-1\normalbaselineskip}[0pt][0pt]{50} & \raisebox{-1\normalbaselineskip}[0pt][0pt]{25} & \raisebox{-1\normalbaselineskip}[0pt][0pt]{87} & \raisebox{-1\normalbaselineskip}[0pt][0pt]{24} & \raisebox{-1\normalbaselineskip}[0pt][0pt]{73} & \raisebox{-1\normalbaselineskip}[0pt][0pt]{10} & \raisebox{-1\normalbaselineskip}[0pt][0pt]{92} & \raisebox{-1\normalbaselineskip}[0pt][0pt]{33} & \raisebox{-1\normalbaselineskip}[0pt][0pt]{56} & \raisebox{-1\normalbaselineskip}[0pt][0pt]{33} & \raisebox{-1\normalbaselineskip}[0pt][0pt]{2} \\
59 & Lucida quae sequi-\newline\hspace*{3mm}tur Canem (α Co-\newline\hspace*{3mm}lumbae). & \raisebox{-2\normalbaselineskip}[0pt][0pt]{35} & \raisebox{-2\normalbaselineskip}[0pt][0pt]{25} & \raisebox{-2\normalbaselineskip}[0pt][0pt]{A} & \raisebox{-2\normalbaselineskip}[0pt][0pt]{18} & \raisebox{-2\normalbaselineskip}[0pt][0pt]{35} & \raisebox{-2\normalbaselineskip}[0pt][0pt]{58} & \raisebox{-2\normalbaselineskip}[0pt][0pt]{55} & \raisebox{-2\normalbaselineskip}[0pt][0pt]{76} & \raisebox{-2\normalbaselineskip}[0pt][0pt]{38} & \raisebox{-2\normalbaselineskip}[0pt][0pt]{120} & \raisebox{-2\normalbaselineskip}[0pt][0pt]{0} & \raisebox{-2\normalbaselineskip}[0pt][0pt]{36} & \raisebox{-2\normalbaselineskip}[0pt][0pt]{34} & \raisebox{-2\normalbaselineskip}[0pt][0pt]{2} \\
60 & Secunda ab ea (β Co-\newline\hspace*{3mm}lumbae). & \raisebox{-1\normalbaselineskip}[0pt][0pt]{37} & \raisebox{-1\normalbaselineskip}[0pt][0pt]{3} & \raisebox{-1\normalbaselineskip}[0pt][0pt]{A} & \raisebox{-1\normalbaselineskip}[0pt][0pt]{16} & \raisebox{-1\normalbaselineskip}[0pt][0pt]{57} & \raisebox{-1\normalbaselineskip}[0pt][0pt]{56} & \raisebox{-1\normalbaselineskip}[0pt][0pt]{44} & \raisebox{-1\normalbaselineskip}[0pt][0pt]{78} & \raisebox{-1\normalbaselineskip}[0pt][0pt]{50} & \raisebox{-1\normalbaselineskip}[0pt][0pt]{123} & \raisebox{-1\normalbaselineskip}[0pt][0pt]{36} & \raisebox{-1\normalbaselineskip}[0pt][0pt]{36} & \raisebox{-1\normalbaselineskip}[0pt][0pt]{39} & \raisebox{-1\normalbaselineskip}[0pt][0pt]{2} \\
61 & [In] rotunditate Na-\newline\hspace*{3mm}vis (ε Navis). & \raisebox{-1\normalbaselineskip}[0pt][0pt]{38} & \raisebox{-1\normalbaselineskip}[0pt][0pt]{30} & \raisebox{-1\normalbaselineskip}[0pt][0pt]{A} & \raisebox{-1\normalbaselineskip}[0pt][0pt]{15} & \raisebox{-1\normalbaselineskip}[0pt][0pt]{30} & \raisebox{-1\normalbaselineskip}[0pt][0pt]{54} & \raisebox{-1\normalbaselineskip}[0pt][0pt]{58} & \raisebox{-1\normalbaselineskip}[0pt][0pt]{124} & \raisebox{-1\normalbaselineskip}[0pt][0pt]{55} & \raisebox{-1\normalbaselineskip}[0pt][0pt]{165} & \raisebox{-1\normalbaselineskip}[0pt][0pt]{23} & \raisebox{-1\normalbaselineskip}[0pt][0pt]{75} & \raisebox{-1\normalbaselineskip}[0pt][0pt]{43} & \raisebox{-1\normalbaselineskip}[0pt][0pt]{2} \\
62 & [In] sectione eius ta-\newline\hspace*{3mm}bulatus (? Navis). & \raisebox{-1\normalbaselineskip}[0pt][0pt]{36} & \raisebox{-1\normalbaselineskip}[0pt][0pt]{19} & \raisebox{-1\normalbaselineskip}[0pt][0pt]{A} & \raisebox{-1\normalbaselineskip}[0pt][0pt]{17} & \raisebox{-1\normalbaselineskip}[0pt][0pt]{41} & \raisebox{-1\normalbaselineskip}[0pt][0pt]{57} & \raisebox{-1\normalbaselineskip}[0pt][0pt]{40} & \raisebox{-1\normalbaselineskip}[0pt][0pt]{128} & \raisebox{-1\normalbaselineskip}[0pt][0pt]{30} & \raisebox{-1\normalbaselineskip}[0pt][0pt]{166} & \raisebox{-1\normalbaselineskip}[0pt][0pt]{10} & \raisebox{-1\normalbaselineskip}[0pt][0pt]{81} & \raisebox{-1\normalbaselineskip}[0pt][0pt]{10} & \raisebox{-1\normalbaselineskip}[0pt][0pt]{2} \\
\cline{2-16}\end{tabular}\end{center}
{\small
\par\hangindent=12mm\hangafter=1 54. Quantitate + 6′ discrepant δ et \textit{h} ab iis quae cum \textit{P} convenirent; in \textit{h} cod. \textarabic{نو} 56′ pro \textarabic{يو} 16′. Iuxta catalogum fixarum deberet \textit{A} esse 70° 53′,6; auctor rotunde 71° posuit, a quibus \textit{B} bene, \textit{C} proxime deducitur. In \textit{B} cod. \textarabic{فه} 85° pro \textarabic{فد} 84°; in \textit{C} vera quantitas 59° 51′.
\par\hangindent=12mm\hangafter=1 55. \textit{P}, δ, \textit{h} bene inter se congruentia; optime \textit{A} cum anno 1211, \textit{B} et \textit{C}.
\par\hangindent=12mm\hangafter=1 56. In \textit{P} cod. \textarabic{نا} 51′ pro \textarabic{يا} 11′. — \textit{P}, δ, \textit{h} inter se convenientes; \textit{A} praecise anno 1211, \textit{B} et \textit{C} respondens.
\par\hangindent=12mm\hangafter=1 57. In \textit{P} cod. \textarabic{قز} 107° pro \textarabic{فز} 87°; \textit{P}, δ, \textit{h} congruentes. At in reliquis supputationibus auctor \textit{P} 87° 30′ posuit, nam hac quantitate et data \textit{A} (optime anno 1211 convenienti) praecise \textit{B} et \textit{C} deducuntur.
\par\hangindent=12mm\hangafter=1 58. In \textit{P} cod. \textarabic{قز} 107° pro \textarabic{فز} 87°; in \textit{C} aliquis minuta rotunda fecit, qua re cod. 30′ pro 33′. \textit{P} cum δ et \textit{h} optime; item \textit{A} cum anno 1211, \textit{B} et \textit{C}.
\par\hangindent=12mm\hangafter=1 59. Cod. in δ \textarabic{ك} 20′ pro \textarabic{كه} 25′; in \textit{h} \textarabic{نح} 58° pro \textarabic{يح} 18°. Bene δ et \textit{h} cum \textit{P}. \textit{A} praecise pro anno 1211: sed miro modo auctor \textit{B} et \textit{C} deduxit \textit{A} 76° 45′ sumendo, quod pro anno 1223 valeret.
\par\hangindent=12mm\hangafter=1 60. Repetitio nr. 16, ubi nomen falsum. In \textit{h} cod. \textarabic{نو} 56° pro \textarabic{يو} 16°.
\par\hangindent=12mm\hangafter=1 61. In \textit{h} cod. \textarabic{نه} 55° pro \textarabic{يه} 15°. Deest in catalogo ob folium deperditum quo stellae Navis continebantur; apud \Name{Ptolemaeum} est 31ª Navis: ὁ ὑποκάτω τῶν γ τῆς ἑπομένης ἀσπιδίσκης « quae est sub tribus « sequentis scutuli ». Iuxta \Name{Baily} ε Navis seu Argûs (Lacaille 830), iuxta \Name{Schjellerup} λ Navis. — A \textit{P} debebant δ 38° 18′,8, \textit{h} 15° 41′,2 deduci; datis \textit{P} et \textit{A} exacte \textit{B} et \textit{C} respondent.
\par\hangindent=12mm\hangafter=1 62. In \textit{h} cod. \textarabic{نز} 57° pro \textarabic{يز} 17°. Apud \Name{Ptolemaeum} 32ª Navis: ὁ ἐπὶ τῆς ἀποτομῆς τοῦ καταστρώματος; iuxta \Name{Baily} Lacaille 864, iuxta \Name{Schjellerup} ψ Navis. — A \textit{P} debebant exactius δ 36° 21′,5, \textit{h} 17° 38′,5 deduci; vel, datis δ et \textit{h}, \textit{P} 57° 43′,4. Inter se \textit{P}, \textit{A}, \textit{B}, \textit{C} optime congruunt.
\par}
\clearpage
\noindent\makebox[\textwidth]{\textbf{}\hfill{\small AL-BATTANI OPUS ASTRONOMICUM}\hfill\textbf{185}}\par\vspace{-1mm}\noindent\rule{\textwidth}{.4pt}\par\vspace{2mm}
\begin{center}\footnotesize\begin{tabular}{r|p{30mm}|r@{\hspace{1.5mm}}r|c|r@{\hspace{1.5mm}}r|r@{\hspace{1.5mm}}r|r@{\hspace{1.5mm}}r|r@{\hspace{1.5mm}}r|r@{\hspace{1.5mm}}r||c|}\cline{2-16}
 & \centering Stellarum nomina. & \multicolumn{2}{c|}{\rotatebox{90}{\parbox{14mm}{\centering\scriptsize Declinatio.}}} & \rotatebox{90}{\parbox{14mm}{\centering\scriptsize Plaga Caeli.}} & \multicolumn{2}{c|}{\rotatebox{90}{\parbox{24mm}{\centering\scriptsize Altitudo meridiana.}}} & \multicolumn{2}{c|}{\rotatebox{90}{\parbox{24mm}{\centering\scriptsize Dimidium commorationis supra Terram.}}} & \multicolumn{2}{c|}{\rotatebox{90}{\parbox{24mm}{\centering\scriptsize Gradus cum quibus culminant.}}} & \multicolumn{2}{c|}{\rotatebox{90}{\parbox{24mm}{\centering\scriptsize Gradus cum quibus oriuntur.}}} & \multicolumn{2}{c|}{\rotatebox{90}{\parbox{24mm}{\centering\scriptsize Gradus cum quibus occidunt.}}} & \rotatebox{90}{\parbox{14mm}{\centering\scriptsize Ordines magnitudinis.}} \\ \hline
63 & Quae est sub tabu-\newline\hspace*{3mm}latu (ζ Navis). & \raisebox{-1\normalbaselineskip}[0pt][0pt]{\textit{37°}} & \raisebox{-1\normalbaselineskip}[0pt][0pt]{\textit{45′}} & \raisebox{-1\normalbaselineskip}[0pt][0pt]{A} & \raisebox{-1\normalbaselineskip}[0pt][0pt]{\textit{59°}} & \raisebox{-1\normalbaselineskip}[0pt][0pt]{\textit{55′}} & \raisebox{-1\normalbaselineskip}[0pt][0pt]{52°} & \raisebox{-1\normalbaselineskip}[0pt][0pt]{14′} & \raisebox{-1\normalbaselineskip}[0pt][0pt]{118°} & \raisebox{-1\normalbaselineskip}[0pt][0pt]{20′} & \raisebox{-1\normalbaselineskip}[0pt][0pt]{162°} & \raisebox{-1\normalbaselineskip}[0pt][0pt]{0′} & \raisebox{-1\normalbaselineskip}[0pt][0pt]{67°} & \raisebox{-1\normalbaselineskip}[0pt][0pt]{40′} & \raisebox{-1\normalbaselineskip}[0pt][0pt]{2} \\
64 & Quae est sub eius\newline\hspace*{3mm}\textit{pectore} (? Navis). & \raisebox{-1\normalbaselineskip}[0pt][0pt]{50} & \raisebox{-1\normalbaselineskip}[0pt][0pt]{42} & \raisebox{-1\normalbaselineskip}[0pt][0pt]{A} & \raisebox{-1\normalbaselineskip}[0pt][0pt]{3} & \raisebox{-1\normalbaselineskip}[0pt][0pt]{18} & \raisebox{-1\normalbaselineskip}[0pt][0pt]{27} & \raisebox{-1\normalbaselineskip}[0pt][0pt]{25} & \raisebox{-1\normalbaselineskip}[0pt][0pt]{112} & \raisebox{-1\normalbaselineskip}[0pt][0pt]{30} & \raisebox{-1\normalbaselineskip}[0pt][0pt]{177} & \raisebox{-1\normalbaselineskip}[0pt][0pt]{24} & \raisebox{-1\normalbaselineskip}[0pt][0pt]{42} & \raisebox{-1\normalbaselineskip}[0pt][0pt]{27} & \raisebox{-1\normalbaselineskip}[0pt][0pt]{2} \\
65 & Praecedens e tribus\newline\hspace*{3mm}eius (? Navis). & \raisebox{-1\normalbaselineskip}[0pt][0pt]{52} & \raisebox{-1\normalbaselineskip}[0pt][0pt]{8} & \raisebox{-1\normalbaselineskip}[0pt][0pt]{A} & \raisebox{-1\normalbaselineskip}[0pt][0pt]{1} & \raisebox{-1\normalbaselineskip}[0pt][0pt]{52} & \raisebox{-1\normalbaselineskip}[0pt][0pt]{20} & \raisebox{-1\normalbaselineskip}[0pt][0pt]{56} & \raisebox{-1\normalbaselineskip}[0pt][0pt]{122} & \raisebox{-1\normalbaselineskip}[0pt][0pt]{43} & \raisebox{-1\normalbaselineskip}[0pt][0pt]{191} & \raisebox{-1\normalbaselineskip}[0pt][0pt]{40} & \raisebox{-1\normalbaselineskip}[0pt][0pt]{45} & \raisebox{-1\normalbaselineskip}[0pt][0pt]{53} & \raisebox{-1\normalbaselineskip}[0pt][0pt]{2} \\
66 & Lucida Hydrae (α Hy-\newline\hspace*{3mm}drae). & \raisebox{-1\normalbaselineskip}[0pt][0pt]{2} & \raisebox{-1\normalbaselineskip}[0pt][0pt]{55} & \raisebox{-1\normalbaselineskip}[0pt][0pt]{A} & \raisebox{-1\normalbaselineskip}[0pt][0pt]{51} & \raisebox{-1\normalbaselineskip}[0pt][0pt]{5} & \raisebox{-1\normalbaselineskip}[0pt][0pt]{87} & \raisebox{-1\normalbaselineskip}[0pt][0pt]{53} & \raisebox{-1\normalbaselineskip}[0pt][0pt]{125} & \raisebox{-1\normalbaselineskip}[0pt][0pt]{58} & \raisebox{-1\normalbaselineskip}[0pt][0pt]{139} & \raisebox{-1\normalbaselineskip}[0pt][0pt]{20} & \raisebox{-1\normalbaselineskip}[0pt][0pt]{107} & \raisebox{-1\normalbaselineskip}[0pt][0pt]{20} & \raisebox{-1\normalbaselineskip}[0pt][0pt]{2} \\
67 & Venter equi Centau-\newline\hspace*{3mm}ri (? Centauri). & \raisebox{-1\normalbaselineskip}[0pt][0pt]{\textit{48}} & \raisebox{-1\normalbaselineskip}[0pt][0pt]{\textit{45}} & \raisebox{-1\normalbaselineskip}[0pt][0pt]{A} & \raisebox{-1\normalbaselineskip}[0pt][0pt]{\textit{5}} & \raisebox{-1\normalbaselineskip}[0pt][0pt]{\textit{15}} & \raisebox{-1\normalbaselineskip}[0pt][0pt]{\textit{35}} & \raisebox{-1\normalbaselineskip}[0pt][0pt]{\textit{10}} & \raisebox{-1\normalbaselineskip}[0pt][0pt]{\textit{192}} & \raisebox{-1\normalbaselineskip}[0pt][0pt]{\textit{30}} & \raisebox{-1\normalbaselineskip}[0pt][0pt]{\textit{239}} & \raisebox{-1\normalbaselineskip}[0pt][0pt]{\textit{0}} & \raisebox{-1\normalbaselineskip}[0pt][0pt]{\textit{118}} & \raisebox{-1\normalbaselineskip}[0pt][0pt]{\textit{10}} & \raisebox{-1\normalbaselineskip}[0pt][0pt]{2} \\
68 & Medium poplitis equi\newline\hspace*{3mm}(ν Centauri). & \raisebox{-1\normalbaselineskip}[0pt][0pt]{53} & \raisebox{-1\normalbaselineskip}[0pt][0pt]{3} & \raisebox{-1\normalbaselineskip}[0pt][0pt]{A} & \raisebox{-1\normalbaselineskip}[0pt][0pt]{0} & \raisebox{-1\normalbaselineskip}[0pt][0pt]{57} & \raisebox{-1\normalbaselineskip}[0pt][0pt]{30} & \raisebox{-1\normalbaselineskip}[0pt][0pt]{0} & \raisebox{-1\normalbaselineskip}[0pt][0pt]{169} & \raisebox{-1\normalbaselineskip}[0pt][0pt]{52} & \raisebox{-1\normalbaselineskip}[0pt][0pt]{221} & \raisebox{-1\normalbaselineskip}[0pt][0pt]{35} & \raisebox{-1\normalbaselineskip}[0pt][0pt]{92} & \raisebox{-1\normalbaselineskip}[0pt][0pt]{0} & \raisebox{-1\normalbaselineskip}[0pt][0pt]{2} \\
69 & Ungula huius equi\newline\hspace*{3mm}(ζ Centauri). & \raisebox{-1\normalbaselineskip}[0pt][0pt]{56} & \raisebox{-1\normalbaselineskip}[0pt][0pt]{37} & \raisebox{-1\normalbaselineskip}[0pt][0pt]{A} & \multicolumn{4}{c|}{\parbox{30mm}{\centering\scriptsize Haec linea est\\semper occulta sub\\horizonte australi,\\sine ortu, qua re}} & \raisebox{-1\normalbaselineskip}[0pt][0pt]{165} & \raisebox{-1\normalbaselineskip}[0pt][0pt]{7} & \multicolumn{4}{c|}{\parbox{30mm}{\centering\scriptsize neque commoratio-\\nem habet supra\\Terram neque cul-\\minationem.}} & \raisebox{-1\normalbaselineskip}[0pt][0pt]{2} \\
70 & Pes huius equi (α Cen-\newline\hspace*{3mm}tauri). & \raisebox{-1\normalbaselineskip}[0pt][0pt]{44} & \raisebox{-1\normalbaselineskip}[0pt][0pt]{38} & \raisebox{-1\normalbaselineskip}[0pt][0pt]{A} & \raisebox{-1\normalbaselineskip}[0pt][0pt]{9} & \raisebox{-1\normalbaselineskip}[0pt][0pt]{22} & \raisebox{-1\normalbaselineskip}[0pt][0pt]{44} & \raisebox{-1\normalbaselineskip}[0pt][0pt]{10} & \raisebox{-1\normalbaselineskip}[0pt][0pt]{177} & \raisebox{-1\normalbaselineskip}[0pt][0pt]{30} & \raisebox{-1\normalbaselineskip}[0pt][0pt]{215} & \raisebox{-1\normalbaselineskip}[0pt][0pt]{46} & \raisebox{-1\normalbaselineskip}[0pt][0pt]{113} & \raisebox{-1\normalbaselineskip}[0pt][0pt]{20} & \raisebox{-1\normalbaselineskip}[0pt][0pt]{1} \\
71 & Genu huius equi (γ\newline\hspace*{3mm}Centauri). & \raisebox{-1\normalbaselineskip}[0pt][0pt]{53} & \raisebox{-1\normalbaselineskip}[0pt][0pt]{57} & \raisebox{-1\normalbaselineskip}[0pt][0pt]{A} & \raisebox{-1\normalbaselineskip}[0pt][0pt]{0} & \raisebox{-1\normalbaselineskip}[0pt][0pt]{3} & \raisebox{-1\normalbaselineskip}[0pt][0pt]{2} & \raisebox{-1\normalbaselineskip}[0pt][0pt]{0} & \raisebox{-1\normalbaselineskip}[0pt][0pt]{191} & \raisebox{-1\normalbaselineskip}[0pt][0pt]{53} & \raisebox{-1\normalbaselineskip}[0pt][0pt]{261} & \raisebox{-1\normalbaselineskip}[0pt][0pt]{30} & \raisebox{-1\normalbaselineskip}[0pt][0pt]{85} & \raisebox{-1\normalbaselineskip}[0pt][0pt]{3} & \raisebox{-1\normalbaselineskip}[0pt][0pt]{2} \\
\cline{2-16}\end{tabular}\end{center}
{\small
\par\hangindent=12mm\hangafter=1 63. In \textit{P} cod. \textarabic{ند} 54′ pro \textarabic{يد} 14′; in \textit{C} \textarabic{سز} 307° pro \textarabic{صز} 67°. Apud \Name{Ptolemaeum} 35ª Navis: ὑπὸ τὸ κατάστρωμα λαμπρός « lucida sub tabulatu »; \Name{Baily} ζ Navis (Lac. 737), \Name{Schjellerup} γ Navis. — \textit{P}, \textit{A}, \textit{B}, \textit{C} optime congruunt; sed δ et \textit{h} falsae, nam ex \textit{P} oriuntur δ 40° 8′, \textit{h} 13° 52′.
\par\hangindent=12mm\hangafter=1 64. In δ cod. \textarabic{نب} 52′ pro \textarabic{مب} 42′; in \textit{h} \textarabic{ح} 8′ pro \textarabic{ىح} = \textarabic{يح} 18′; in \textit{P} \textarabic{\AbjadZero{}} 0′ pro \textarabic{كه} 25′. His correctis, omnia bene congruunt. — Apud \Name{Ptol.} 36ª Navis: ἐπὶ τῆς κάτω τρόπεως λαμπρός « lucida infra, in carina »; \Name{Baily} η Navis (Lac. 721), \Name{Schjellerup} χ Navis.
\par\hangindent=12mm\hangafter=1 65. In δ cod. \textarabic{ب} 2′ pro \textarabic{ح} 8′; in \textit{A} \textarabic{قلب} 132° pro \textarabic{قكب} 122°; quibus correctis omnia optime congruunt. — Apud \Name{Ptol.} 37ª Navis: τῶν ἑπομένων τούτῳ γ ὁ προηγούμενος « anterior trium hanc sequentium »; \Name{Baily} \textit{q} Navis (Lac. 786), \Name{Schjellerup} δ Navis.
\par\hangindent=12mm\hangafter=1 66. In \textit{P} cod. \textarabic{يج} 13′ pro \textarabic{نج} 53′. Omnia optime. In catalogo, ob folium deperditum, deest ut nr. 61-65.
\par\hangindent=12mm\hangafter=1 67. Nullo modo elementa huius stellae \Name{Schiaparelli} emendare potuit. Fortasse δ et \textit{A} potius 23ª stellarum Centauri in catalogo (p. 174) indicarent; sed reliqui numeri non concordant.
\par\hangindent=12mm\hangafter=1 68. In Catalogo est 20ª stella Centauri: « in poplite dextero prope pedem » (p. 174). In δ cod. \textarabic{يج} 13° pro \textarabic{نج} 53°. Optime congruunt inter se \textit{h} et δ; item optime inter se \textit{P}, \textit{A}, \textit{B}, \textit{C}; sed inter δ (et \textit{h}) et \textit{P} magna habetur discrepantia. Error ex eo provenit quod auctor, in \textit{B} et \textit{C} supputandis, arcum diurnum pro arcu semidiurno sumpsit; i. e. posuit \textit{P} 30° 0′ (ut cod. habet) loco 15° 0′. \textit{A} anno 1211 non convenit. — Iuxta \Name{Baily} ν Centauri (Lac. 1070); \Name{Schjellerup} γ Crucis.
\par\hangindent=12mm\hangafter=1 69. In δ cod. \textarabic{يو} 16° pro \textarabic{نو} 56° (iuxta catalogum 57° 17′ pro 1211). \textit{A} est 4′ maior quam quae anno 1211 conveniret. Iuxta \Name{Baily} ζ Centauri (Lac. 1082); \Name{Schjellerup} α Crucis. In urbe ar-Raqqah non aspicitur. Magnitudo ut S. et Pts.; Almagesta graeca 4.
\par\hangindent=12mm\hangafter=1 70. In \textit{P} cod. \textarabic{نج} 53′ pro \textarabic{ى} 10′; error ab auctore provenit, qui, in 90° − \textit{m} = \textit{P} scribendo, legit \textit{m} 45° 7′ loco 45° 50′ (\textarabic{ز} pro \textarabic{ن}). Sed omnia optime congruunt cum \textit{P} 44° 10′, quem recepi. \textit{A} est 6′,6 maior quam quae anno 1211 conveniret.
\par\hangindent=12mm\hangafter=1 71. In \textit{B} cod. \textarabic{\AbjadZero{}} 0′ pro \textarabic{ل} 30′; in \textit{C} \textarabic{قكه} 125° pro \textarabic{فه} 85°. Quantitate − 2′ discrepat δ ab ea quam \textit{P} praeberet. Datis \textit{P} et \textit{A} (quae tamen + 1° 20′ ab epocha 1211 differt) deducuntur satis bene \textit{B} et \textit{C} (ve-
\par}
\par\vfill\hfill 24
\clearpage
\noindent\makebox[\textwidth]{\textbf{186}\hfill{\small C. A. NALLINO}\hfill\textbf{}}\par\vspace{-1mm}\noindent\rule{\textwidth}{.4pt}\par\vspace{2mm}
\begin{center}\footnotesize\begin{tabular}{r|p{30mm}|r@{\hspace{1.5mm}}r|c|r@{\hspace{1.5mm}}r|r@{\hspace{1.5mm}}r|r@{\hspace{1.5mm}}r|r@{\hspace{1.5mm}}r|r@{\hspace{1.5mm}}r||c|}\cline{2-16}
 & \centering Stellarum nomina. & \multicolumn{2}{c|}{\rotatebox{90}{\parbox{14mm}{\centering\scriptsize Declinatio.}}} & \rotatebox{90}{\parbox{14mm}{\centering\scriptsize Plaga Caeli.}} & \multicolumn{2}{c|}{\rotatebox{90}{\parbox{24mm}{\centering\scriptsize Altitudo meridiana.}}} & \multicolumn{2}{c|}{\rotatebox{90}{\parbox{24mm}{\centering\scriptsize Dimidium commorationis supra Terram.}}} & \multicolumn{2}{c|}{\rotatebox{90}{\parbox{24mm}{\centering\scriptsize Gradus cum quibus culminant.}}} & \multicolumn{2}{c|}{\rotatebox{90}{\parbox{24mm}{\centering\scriptsize Gradus cum quibus oriuntur.}}} & \multicolumn{2}{c|}{\rotatebox{90}{\parbox{24mm}{\centering\scriptsize Gradus cum quibus occidunt.}}} & \rotatebox{90}{\parbox{14mm}{\centering\scriptsize Ordines magnitudinis.}} \\ \hline
72 & Lucida e stellis Co-\newline\hspace*{3mm}ronae australis (θ\newline\hspace*{3mm}Cor. Austr.). & \raisebox{-2\normalbaselineskip}[0pt][0pt]{40°} & \raisebox{-2\normalbaselineskip}[0pt][0pt]{44′} & \raisebox{-2\normalbaselineskip}[0pt][0pt]{A} & \raisebox{-2\normalbaselineskip}[0pt][0pt]{13°} & \raisebox{-2\normalbaselineskip}[0pt][0pt]{16′} & \raisebox{-2\normalbaselineskip}[0pt][0pt]{51°} & \raisebox{-2\normalbaselineskip}[0pt][0pt]{14′} & \raisebox{-2\normalbaselineskip}[0pt][0pt]{268°} & \raisebox{-2\normalbaselineskip}[0pt][0pt]{27′} & \raisebox{-2\normalbaselineskip}[0pt][0pt]{288°} & \raisebox{-2\normalbaselineskip}[0pt][0pt]{10′} & \raisebox{-2\normalbaselineskip}[0pt][0pt]{248°} & \raisebox{-2\normalbaselineskip}[0pt][0pt]{11′} & \raisebox{-2\normalbaselineskip}[0pt][0pt]{4} \\
73 & Initium cursus a-\newline\hspace*{3mm}quae ab ansa am-\newline\hspace*{3mm}phorae (? Aquarii). & \raisebox{-2\normalbaselineskip}[0pt][0pt]{16} & \raisebox{-2\normalbaselineskip}[0pt][0pt]{21} & \raisebox{-2\normalbaselineskip}[0pt][0pt]{A} & \raisebox{-2\normalbaselineskip}[0pt][0pt]{37} & \raisebox{-2\normalbaselineskip}[0pt][0pt]{39} & \raisebox{-2\normalbaselineskip}[0pt][0pt]{77} & \raisebox{-2\normalbaselineskip}[0pt][0pt]{39} & \raisebox{-2\normalbaselineskip}[0pt][0pt]{323} & \raisebox{-2\normalbaselineskip}[0pt][0pt]{0} & \raisebox{-2\normalbaselineskip}[0pt][0pt]{325} & \raisebox{-2\normalbaselineskip}[0pt][0pt]{45} & \raisebox{-2\normalbaselineskip}[0pt][0pt]{321} & \raisebox{-2\normalbaselineskip}[0pt][0pt]{30} & \raisebox{-2\normalbaselineskip}[0pt][0pt]{} \\
74 & Medium cursus a-\newline\hspace*{3mm}quae ab ansa (? A-\newline\hspace*{3mm}quarii). & \raisebox{-2\normalbaselineskip}[0pt][0pt]{10} & \raisebox{-2\normalbaselineskip}[0pt][0pt]{7} & \raisebox{-2\normalbaselineskip}[0pt][0pt]{A} & \raisebox{-2\normalbaselineskip}[0pt][0pt]{43} & \raisebox{-2\normalbaselineskip}[0pt][0pt]{53} & \raisebox{-2\normalbaselineskip}[0pt][0pt]{82} & \raisebox{-2\normalbaselineskip}[0pt][0pt]{33} & \raisebox{-2\normalbaselineskip}[0pt][0pt]{325} & \raisebox{-2\normalbaselineskip}[0pt][0pt]{46} & \raisebox{-2\normalbaselineskip}[0pt][0pt]{322} & \raisebox{-2\normalbaselineskip}[0pt][0pt]{45} & \raisebox{-2\normalbaselineskip}[0pt][0pt]{327} & \raisebox{-2\normalbaselineskip}[0pt][0pt]{30} & \raisebox{-2\normalbaselineskip}[0pt][0pt]{4} \\
75 & Finis cursus aquae\newline\hspace*{3mm}ab ansa (χ Aqua-\newline\hspace*{3mm}rii). & \raisebox{-2\normalbaselineskip}[0pt][0pt]{12} & \raisebox{-2\normalbaselineskip}[0pt][0pt]{25} & \raisebox{-2\normalbaselineskip}[0pt][0pt]{A} & \raisebox{-2\normalbaselineskip}[0pt][0pt]{41} & \raisebox{-2\normalbaselineskip}[0pt][0pt]{35} & \raisebox{-2\normalbaselineskip}[0pt][0pt]{80} & \raisebox{-2\normalbaselineskip}[0pt][0pt]{48} & \raisebox{-2\normalbaselineskip}[0pt][0pt]{332} & \raisebox{-2\normalbaselineskip}[0pt][0pt]{26} & \raisebox{-2\normalbaselineskip}[0pt][0pt]{333} & \raisebox{-2\normalbaselineskip}[0pt][0pt]{0} & \raisebox{-2\normalbaselineskip}[0pt][0pt]{330} & \raisebox{-2\normalbaselineskip}[0pt][0pt]{33} & \raisebox{-2\normalbaselineskip}[0pt][0pt]{4} \\
\cline{2-16}\end{tabular}\end{center}
{\small
\par\hangindent=12mm\hangafter=1 rae quantitates 261° 26′,6, 84° 58′,1). — Iuxta \Name{Baily} γ Centauri (Lac. 1185); \Name{Schjellerup} β Centauri.
\par\hangindent=12mm\hangafter=1 72. In \textit{h} cod. \textarabic{لو} 36′ pro \textarabic{ىو} = \textarabic{يو} 16′; in \textit{P} \textarabic{يا ند} 11° 54′ pro \textarabic{نا يد} 51° 14′; in \textit{C} \textarabic{د} 4′ pro \textarabic{يا} 11′. Quibus correctis, omnia optime congruunt; sed \textit{A} discrepat + 13′,1 ab epocha 1211, qua re potius ad 1223 pertinet.
\par\hangindent=12mm\hangafter=1 73. Omnia satis bene inter se congruentia: sed ad locum constellationis Aquarii pertinent (long. 322° 12′ pro 1211, vel 321° 54′ pro anno 1191 catalogi; lat. austr. 2° 17′), in quo nullum sidus catalogus Albatenianus indicat.
\par\hangindent=12mm\hangafter=1 74. In \textit{P} cod. \textarabic{قيا نج} 111° 53′ pro \textarabic{فب لج} 82° 33′; in \textit{C} \textarabic{سلز} 337° pro \textarabic{سكز} 327°. Optime δ et \textit{h} cum \textit{P}; exacte \textit{A} pro anno 1211; satis proxime \textit{B} et \textit{C}. — Est tertiadecima stella Aquarii in catalogo, p. 165.
\par\hangindent=12mm\hangafter=1 75. In \textit{P} cod. \textarabic{قب} 102° pro \textarabic{ف} 80°; in \textit{B} \textarabic{سلب ك} 332° 20′ pro \textarabic{سلج لا} 333° 0′. Omnia optime; \textit{A} praecise pro anno 1211. — Est septimadecima stellarum Aquarii, p. 165.
\par}
\end{document}
